\section{The window bound: hypotheses H and E}\label{sec:pf:vwin-he}

\providecommand{\vwinheQZ}{\Q_\ell/\Z_\ell}
\providecommand{\vwinhecyc}{\chi_{\mathrm{cyc}}}
\providecommand{\vwinheVt}{\widetilde V}
\providecommand{\vwinheX}{\mathcal X}
\providecommand{\vwinheA}{\mathcal A}
\providecommand{\vwinhesat}{\operatorname{sat}}
\providecommand{\vwinhecot}{\operatorname{Cot}}
\providecommand{\vwinhecotlen}{\lambda^{\mathrm{cot}}}
\providecommand{\vwinheloclen}{\lambda}
\providecommand{\vwinhesecs}{\operatorname{secs}}
\providecommand{\vwinheIrr}{\operatorname{Irr}}
\providecommand{\vwinheAnn}{\operatorname{Ann}}
\providecommand{\vwinheD}{\mathfrak D}
\providecommand{\vwinheXX}{\mathfrak X}
\providecommand{\vwinhecoker}{\lambda^{\mathrm{co}}}

This section proves \cref{thm:window:HE}: the package $\sfname{HypothesesH'}$ of theorems of the formal
development implies the window bound $\sfname{VWin}$ of \cref{def:window:vwin}. The proof
(\cref{sec:pf:vwin-he:proof}) is a short chain of nine results, each proved elsewhere in Part II. From
$\sfname{HypothesesH'}$ they produce Lemma K (\cref{thm:window:lemmaK}), the vanishing theorem for Bloch--Kato
Selmer groups (\cref{thm:kato}), automorphy statements for three- and four-dimensional Galois representations,
the counting statements $\sfname{CountI}$, $\sfname{CountII}$, $\sfname{CountIII}$, $\sfname{CountV}$ for the plus
part, and finally $\sfname{VWin}$. The counting statements are made precise in \cref{sec:pf:vwin-he:counts}; the
package and the intermediate statements are described in \cref{sec:pf:vwin-he:hyps}.

The remaining subsections treat two auxiliary results which enter the proofs of several of these inputs.
\Cref{sec:pf:vwin-he:lemmaM} is devoted to Lemma M, a bound for the length of the cotangent module in degree $-1$
of a $\Z$-graded algebra at a fixed point $x$ of the grading torus, modulo given equations, by twice the sum of the
valuations of a degree-zero function at representatives of the orbits of the other points. It is the commutative
algebra behind step (ii) of the Eisenstein congruence method of \cref{sec:window:eis}, and it is used in the proofs
of $\sfname{CountI}$, $\sfname{CountII}$, $\sfname{CountIII}$ and $\sfname{CountV}$. Its proof compares two
estimates for Euler characteristics of Koszul type on a trigraded Rees algebra: a lower bound computed on a formal
chart at $x$, and an upper bound computed from the points. We give this argument completely; the auxiliary lemmas
are stated where they are used and proved in \cref{sec:pf:vwin-he:auxproofs}, and a few deeper ingredients
(finiteness of the points on a slice, the prime ideals of the Rees algebra, a separation property of charts, and
three further estimates) are cited with precise statements. \Cref{sec:pf:vwin-he:descent} proves the descent formula along
cyclotomic $\Z_\ell$-towers underlying \cref{lem:window:tower}: the length of the Tate--Shafarevich group of a
Tate twist over a layer of the cyclotomic tower equals the length of the cokernel of $\gamma^{\ell^m}-1$ on a
twisted Iwasawa module. This proof is complete.

\subsection{Coefficient rings, hosts and the counting statements}\label{sec:pf:vwin-he:counts}

Throughout this section $\Omega$ is an algebraic closure of $\Q$ and $G_\Q=\Gal(\Omega/\Q)$, with the Krull
topology. For a prime $\ell$ we fix an algebraic closure $\Qlbar$ of $\Ql$, and $\vwinhecyc:G_\Q\to\Zl^\times$ is
the $\ell$-adic cyclotomic character. Cohomology of profinite groups is continuous cohomology; modules of the form
$\OO\otimes\vwinheQZ$ carry the discrete topology. A valuation ring of $\Omega$ is any valuation subring
$A\subset\Omega$ (the trivial one $A=\Omega$ included); $D_A=\{g\in G_\Q: gA=A\}$ and
$I_A=\{g\in D_A: v_A(ga-a)<1\ \text{for all}\ a\in A\}$ are its decomposition and inertia groups. For the trivial
valuation ring, $I_\Omega=1$.

\begin{definition}[Coefficient rings and normalized lengths]\label{pf:vwin-he:coef}
Let $\ell$ be a prime. A \emph{coefficient ring} for $\ell$ is a $\Zl$-subalgebra $\OO\subset\Qlbar$ which is a
discrete valuation ring and a finitely generated $\Zl$-module. Put $e_\OO=\len_\OO(\OO/\ell\OO)$, the
ramification index of $\OO$ over $\Zl$. For $n\in\N\cup\{\infty\}$ put $\|n\|_\OO=n/e_\OO$ if $n<\infty$, and
$\|\infty\|_\OO=0$.
\end{definition}

The convention $\|\infty\|_\OO=0$ is that of the formal development, where a length in $\N\cup\{\infty\}$ enters
real inequalities through its natural-number part. For this reason every bound below for normalized lengths is
accompanied by a separate finiteness assertion.

\begin{definition}[The lengths $V_\ell(\theta,j)$]\label{pf:vwin-he:V}
Let $\ell$ be a prime, $\OO$ a coefficient ring for $\ell$, $N\ge1$, $\theta$ a Dirichlet character modulo $N$
with values in $\OO$, and $j\ge0$. Let $\kappa_N:G_\Q\to(\Z/N)^\times$ be the cyclotomic character modulo $N$
($g\zeta=\zeta^{\kappa_N(g)}$ for $\zeta^N=1$), and $\theta_G=\theta\circ\kappa_N:G_\Q\to\OO^\times$. Let
$M_{\theta,j}=\OO\otimes_{\Zl}\vwinheQZ$, on which $G_\Q$ acts through $\theta_G^{-1}\vwinhecyc^{-j}$. For $n\ge1$
let $\Sha_n(\Q,M_{\theta,j})\subset H^1(G_\Q,M_{\theta,j})$ be the submodule of classes whose restrictions to
$I_A$ for every valuation ring $A$ of $\Omega$, and to $D_A$ for every valuation ring $A$ whose maximal ideal
contains $n$, vanish. Put
\[
 V_\ell(\theta,j)=\len_\OO\Sha_{\ell N}(\Q,M_{\theta,j})\in\N\cup\{\infty\},\qquad
 \vwinheVt_\ell(\theta,j)=\|V_\ell(\theta,j)\|_\OO .
\]
For $B\in\R$ we say that $(\theta,j)$ is \emph{bounded by $B$} if $V_\ell(\theta,j)<\infty$ and
$\vwinheVt_\ell(\theta,j)\log\ell\le B$.
\formallabel{UpperTheory.V.HV.VLe}
\end{definition}

These are the lengths of \cref{def:window:host}; they depend on $\Omega$, which is fixed in each statement below.

\begin{definition}[Hosts]\label{pf:vwin-he:host}
A \emph{host} $h=(\chi,m,q_1,d)$ consists of a primitive Dirichlet character $\chi$ with complex values, of
conductor $N_\chi$ and order $o_\chi$, an integer $m\ge2$, and primes $q_1\equiv3\pmod8$, $d\equiv7\pmod{16}$ with
$q_1d$ prime to $N_\chi$. Put $N_h=\lcm(2^{m+2},N_\chi)\,q_1d$ and let $H_h\subset(\Z/N_h)^\times$ be the
subgroup of classes $a$ with $a\equiv\pm1\pmod{2^{m+2}}$, $\chi(a)=1$, and $a$ a square modulo $q_1$ and modulo
$d$. For a prime $\ell$, a coefficient ring $\OO$ for $\ell$ and $j\ge0$, let $\vwinheX_h(\OO,j)$ be the finite
set of Dirichlet characters $\theta$ modulo $N_h$ with values in $\OO$ such that $H_h\subset\ker\theta$ and
$\theta(-1)=(-1)^j$. With $p_K(\chi)$, $h^*(\chi)$, $p_1(\chi)$ and the bad set $B_\chi$ of
\cref{def:window:cases}, the host $h$ is \emph{big for} $t\in\N$ if $2t+1<\min(q_1,d)$ and each of $N_\chi$,
$o_\chi$, $p_K(\chi)$, $h^*(\chi)$, $p_1(\chi)$ is smaller than $\min(q_1,d)$.
\end{definition}

\begin{definition}[Components of characters, and the cases]\label{pf:vwin-he:cases}
Let $N\ge1$ and let $\psi$ be a Dirichlet character modulo $N$ with values in a commutative ring. For a prime $r$
let $\psi_{(r)}=\psi\circ\pi_r$, where $\pi_r(u)\in(\Z/N)^\times$ is congruent to $u$ modulo the $r$-part of $N$
and to $1$ modulo its prime-to-$r$ part. For a prime $\ell$ write $\psi=\psi^{(\ell)}\psi_{\ell'}$ with
$\psi^{(\ell)}$ of $\ell$-power order and $\psi_{\ell'}$ of order prime to $\ell$, and put
\[
 \psi^\flat=\psi_{\ell'}\,(\psi_{\ell'})_{(\ell)}^{-1},\qquad \psi^\sharp=\psi_{\ell'}\,(\psi_{\ell'})_{(2)}^{-1}.
\]
Let $f_0(\ell,\psi)$ be the order of $\ell$ modulo $\lcm(4,\ord\psi^\sharp)$, $n_0(\ell,\psi)=v_2(\ell^{f_0(\ell,\psi)}-1)$,
and $c_2(\psi)=v_2(\cond\psi_{(2)})$. For a host $h=(\chi,m,q_1,d)$, a prime $\ell$, $j\ge1$ and a character
$\theta$ modulo $N_h$, put $R_1(h,\ell,\theta)=16\,o_\chi(\ell N_\chi)^2\,4^{n_0(\ell,\theta)}$, and say that
$(\ell,j,\theta)$ is
\begin{enumerate}[label=(\Roman*)]
\item[(I)] in case I if $\ell\notin B_\chi$ and $(\ell-1)\nmid2j$;
\item[(II)] in case II if $j\ge2$, $\ord\theta^\flat$ is even, and $\ell\in B_\chi$ or $(\ell-1)\mid2j$;
\item[(III)] in case III if $\ell\notin B_\chi$, $(\ell-1)\mid 2j$, $\ord\theta$ is odd and $\theta\ne1$;
\item[(V)] in case V if $\ell\in B_\chi$ and $j=1$.
\end{enumerate}
\formallabel{UpperTheory.V.HV.n0}
\end{definition}

The character $\theta^\flat$ is the character of \cref{def:window:cases}, and the four cases are those of
\cref{tab:window:cases}.

\begin{definition}[Counting statements]\label{pf:vwin-he:countdef}
For a property $P$ of tuples $(h,\ell,\OO,j,\theta)$, let $\mathrm{Count}(P)$ be the statement: \emph{for every
primitive Dirichlet character $\chi$ there is $C\in\N$ such that for every host $h$ with character $\chi$, every
$t\in\N$ for which $h$ is big, every algebraic closure $\Omega$ of $\Q$, every odd prime $\ell$, every coefficient
ring $\OO$ for $\ell$, every $0<j<t$ and every $\theta\in\vwinheX_h(\OO,j)$ with $P(h,\ell,\OO,j,\theta)$, the pair
$(\theta,j)$ is bounded by $(\ell N_hj)^C$.} Then
\begin{itemize}
\item $\sfname{CountII}$ and $\sfname{CountIII}$ are $\mathrm{Count}(P)$ for $P$: ``$(\ell,j,\theta)$ is in
case II'', resp.\ ``in case III'';
\item $\sfname{CountV}$ is $\mathrm{Count}(P)$ for $P$: ``$(\ell,j,\theta)$ is in case V and
$R_1(h,\ell,\theta)\le2^{c_2(\theta)}$'';
\item $\sfname{CountI}$ is the statement: \emph{for every primitive $\chi$ there are $C\in\N$ and primes $r_1,r_2$
such that for every host $h=(\chi,m,q_1,d)$ with $r_1,r_2<\min(q_1,d)$, every $t$ for which $h$ is big, every
$0<j<t$, every $\Omega$, every choice of a coefficient ring $\OO_\ell$ for each prime $\ell$, and every finite set
$L$ of odd primes $\ell$ such that $(\ell,j)$ is in case I, one has $V_\ell(\theta,j)<\infty$ for $\ell\in L$ and
$\theta\in\vwinheX_h(\OO_\ell,j)$, and}
\[
 \sum_{\ell\in L}\log\ell\sum_{\theta\in\vwinheX_h(\OO_\ell,j)}\vwinheVt_\ell(\theta,j)\ \le\ (N_ht)^C .
\]
\end{itemize}
\formallabel{UpperTheory.V.HV.CountShape}
\formallabel{UpperTheory.V.HV.CountI}
\formallabel{UpperTheory.V.HV.CountII}
\formallabel{UpperTheory.V.HV.CountIII}
\formallabel{UpperTheory.V.HV.CountV}
\end{definition}

These are the statements of \cref{def:window:counts}; the auxiliary $2$-adic condition of case V is
$R_1(h,\ell,\theta)\le2^{c_2(\theta)}$. Case I is the only case summed over primes; in the other cases the bound
depends on $\ell$.

\subsection{The package H and the intermediate statements}\label{sec:pf:vwin-he:hyps}

The hypothesis of \cref{thm:window:HE} is a conjunction of statements, each of which is a theorem of the formal
development \cite{LeanSource2026}; most are formalized published theorems, a few are formal proofs of known
theorems (see \cref{sec:window:arch} and \cref{rem:window:bc}). Its only role in this section is to be passed to
the results cited in \cref{sec:pf:vwin-he:proof}, so we list its components briefly.

\begin{definition}[The package $\sfname{HypothesesH'}$]\label{pf:vwin-he:H}
$\sfname{HypothesesH'}$ is the conjunction of the following statements.
\begin{enumerate}[label=(H\arabic*)]
\item\label{pf:vwin-he:H1} $\sfname{LiteratureRankNHypothesis}$, which consists of four groups of statements:
 (a) the representation theory of $\GL_n$ over $p$-adic fields used in the development (finiteness theorems of
 Bernstein--Zelevinsky, the geometric lemma, Zelevinsky's classification and polynomial ring, the Langlands
 classification by multisegments, standard modules and their cosocles, irreducibility criteria for parabolic
 induction, square-integrable and tempered representations, the Moy--Prasad isomorphism, depth-zero induction,
 the local converse theorem, Brauer's theorem in degree zero, Chevalley's unit theorem and Green's regular
 characters) \cite{BernsteinZelevinsky1977,Zelevinsky1980,MoyPrasad1994};
 (b) automorphic forms on $\GL_n$ and Galois representations: rapid decay of cusp forms, Harish-Chandra's
 finiteness theorem, finiteness of automorphic forms of fixed type, the theorem of Jacquet--Shalika and the
 cuspidal spectrum \cite{JacquetShalika1981}, base change from definite unitary groups to isobaric representations
 (\cref{rem:window:bc}(i)), Gelfand--Kazhdan duality, the existence of Galois representations attached to
 polarized regular algebraic cuspidal representations, of the type of \cref{thm:row0b}, Clozel's purity
 and cohomologicality theorems \cite{Clozel1990}, global Artin reciprocity and ray class fields, $\ell$-adic
 characters attached to algebraic Hecke characters, and Berger's theorem that de Rham representations are
 potentially semistable \cite{Berger2002};
 (c) Galois cohomology: Poitou--Tate duality, finiteness, Tate duality and Euler characteristics for local fields
 (also for finite extensions of $\Ql$), global finiteness and Euler characteristics, orthogonality of unramified
 classes, the Greenberg--Wiles formula \cite{Wiles1995} and Bloch--Kato's local theory \cite{BlochKato1990};
 (d) complex multiplication: Grothendieck's deformation theory, de Rham cohomology, dual crystals, period
 comparison and Gauss--Manin periods of abelian varieties, the Serre--Tate good reduction criterion, the main
 theorem of complex multiplication for elliptic curves, the Hilbert class field via $j$, the theorem of
 Siegel--Klingen--Deligne--Ribet \cite{DeligneRibet1980}, and the functional equations of partial zeta functions
 and Hecke $L$-functions.
\item\label{pf:vwin-he:H2} Further statements on automorphic forms: Clozel's criterion for cohomological cuspidal
 representations of $\GL_a$ over totally imaginary fields and the existence of $\Aut(\C)$-conjugates of regular
 algebraic cuspidal representations \cite{Clozel1990}; Borel's bound for the number of cohomological cuspidal
 representations of given level by the dimension of the cohomology of the locally symmetric space
 \cite{Borel1983}; Galois representations with local--global compatibility for regular algebraic polarized cuspidal
 representations of $\GL_2$ over totally real fields \cite{Carayol1986,Taylor1989}; strong multiplicity one,
 Gelfand--Kazhdan duality (the contragredient is the transpose-inverse) and admissibility for cuspidal
 representations of $\GL_a$ over totally real fields; quadratic base change for $\GL_2$ from a totally real field to
 a totally imaginary quadratic extension \cite{Langlands1980}; and the cuspidal representation of $\GL_2(\A_\Q)$
 attached to a classical newform of weight $\ge2$.
\item\label{pf:vwin-he:H3} The base change and descent statements of \cref{rem:window:bc}: descent of regular
 algebraic conjugate self-dual cuspidal representations $\pi$ of $\GL_n(\A_E)$, $2\le n\le4$, $E$ a CM field
 unramified over $E^+$ at all finite places, to an automorphic representation of the definite unitary group which is
 unramified at the non-split places and whose components at almost all split places have the Satake parameters of
 $\pi$ (cf.\ \cite{Labesse2011}); cyclic base change and descent of prime degree over totally imaginary fields
 \cite{ArthurClozel1989}; descent to $\Q$ of regular algebraic cuspidal representations of $\GL_2$ over an imaginary
 quadratic field with sufficiently regular infinitesimal character and Galois-invariant Satake parameters at almost
 all places, or an alternative Galois-theoretic form of this descent; the existence of a totally definite quaternion
 algebra over every totally real field of even degree \cite{Vigneras1980}; and the semisimplicity of the space of
 cusp forms over totally real fields.
\end{enumerate}
\formal{UpperTheory.V.Af.HypothesesH'}
\end{definition}

The intermediate statements of the proof are the following. Two of them are recalled from Part I:
$\sfname{LemmaK}(C_K)$, for $C_K\in\R$, is the bound of \cref{def:window:lemmaK}, namely
$\log\#H^2(G_{L,S_\ell},\Zl(j+1))\le\exp\bigl((2j[L:\Q])^{C_K}\log(2|d_L|)\bigr)$ for all totally imaginary $L$
containing $\mu_\ell$ ($\mu_4$ if $\ell=2$) and all $j\ge1$; and $\sfname{KatoTwoW}$ is the vanishing of the
Bloch--Kato Selmer groups of $W$ and $W\otimes\vartheta$ for the two-dimensional representations $W$ of
\cref{def:window:kato}. The other four are automorphy statements for the deformation problems used for the plus
part. They occur below only as conclusions and hypotheses of cited theorems, so we only describe their content.

\begin{definition}[Automorphy statements for the plus part]\label{pf:vwin-he:autstmts}
\leavevmode
\begin{enumerate}[label=(\arabic*)]
\item\label{pf:vwin-he:Sp} \emph{Symplectic automorphy} $(\mathrm{Sp})$: for every host $h$ big for $t$, every odd
prime $\ell$, every coefficient ring $\OO$, every $0<j<t$, every $\theta\in\vwinheX_h(\OO,j)$ and every admissible
four-dimensional deformation datum of $\GSp_4$-type attached to $(\ell,j,\theta)$ in case II or III, the following
hold: the weight-zero subalgebras generated by the points of the datum are finite $\OO$-modules; at the
block-diagonal characteristic-zero points, the values of the upper $2\times2$ block span a proper subspace of the
$2\times2$ matrices; and there is an isomorphism $\iota:\Qlbar\cong\C$ such that every irreducible
characteristic-zero point $r$ of the datum is automorphic in the following sense: for every imaginary quadratic field
$K$ with $r|_{G_K}$ irreducible, $r|_{G_K}$ is attached via $\iota$ to a regular algebraic cuspidal representation of
$\GL_4(\A_K)$, and a window condition on Satake parameters holds.
\item\label{pf:vwin-he:Ue} \emph{Unitary automorphy at the ordinary points} $(\mathrm{U}_{\mathrm{ord}})$: for every
prime $\ell$, every coefficient ring $\OO$, every $j$ and Dirichlet character $\theta$ with values in $\OO$, every
datum $D$ consisting of an imaginary quadratic field $K$ (with quadratic character $\eta$) satisfying the genericity
conditions of case I for $(K,\ell,j,\theta)$, a complex conjugation and an $\ell$-adic place, and every place $B$ of
$\Omega$ above $\ell$, there is an element $g_1\in G_\Q$ with the marking property such that: (1) on some subgroup of
finite index of $\ker\eta$, the characteristic polynomials of all ordinary three-dimensional points of unitary type
attached to $(\eta,\Psi=\vwinhecyc^j\theta_G,g_1,B)$ with coefficients in algebras over the residue field of $\OO$
agree with those of the base representation; (2) every such ordinary point with coefficients in a finite discrete
valuation ring inside $\Qlbar$ which is irreducible on every open subgroup of $\ker\eta$ becomes, after restriction
to $G_K$ and twist by $\vwinhecyc^{-1}$, automorphic of conjugate self-dual type over $K$; and (3) at the endoscopic
such points, every two-dimensional representation $W$ agreeing with the two-dimensional block on $\ker\eta$ and of
trace $0$ at the complex conjugation is such that $W\otimes\vwinhecyc^{-j}$ is attached to a newform of weight
$2j+1$, and $W$ satisfies the window condition $\sfname{WitAQ}$ of \cref{def:window:kato}.
\item\label{pf:vwin-he:Ui} \emph{Unitary automorphy over $\Q(i)$} $(\mathrm{U}_{i})$: for $K=\Q(i)$ and all host
data, the universal ring of every host frame (the three-dimensional unitary deformation problem used in case V)
has a weight-zero part which is a finite $\OO$-module, and all its irreducible characteristic-zero points are
automorphic over $K$.
\item\label{pf:vwin-he:Uu} \emph{Unitary automorphy with unramified coefficients} $(\mathrm{U}_{\mathrm{ur}})$: for
every datum $D$ as in $(\mathrm{U}_{\mathrm{ord}})$ with unramified coefficient ring ($e_\OO=1$) and every
$\iota:\Qlbar\cong\C$, there is a marked Frobenius element such that the weight-zero subalgebras of the relevant
deformation rings are finite over $\OO$, their irreducible characteristic-zero points are automorphic, and at their
endoscopic points the two-dimensional block is modular in the sense used in case I.
\end{enumerate}
\formal{UpperTheory.V.Aut.SymplecticFourImageDense}
\formal{UpperTheory.V.Pr.GdEndAlt}
\formal{UpperTheory.V.Aut.ThmS6Core_s}
\formal{UpperTheory.V.Aut.AutomorphyUnramifiedCoefficientsAlt}
\end{definition}

Statement $(\mathrm{U}_{\mathrm{ur}})$ is the ``automorphy statement of $R=T$ type'' of \cref{thm:window:countI}.

\subsection{Proof of the main reduction}\label{sec:pf:vwin-he:proof}

\begin{theorem}[\cref{thm:window:HE}]\label{pf:vwin-he:main}
$\sfname{HypothesesH'}$ implies $\sfname{VWin}$.
\formallabel{UpperTheory.V.HV.DescentVanishingReductions.vwin_of_hypotheses_H_E_alt}
\end{theorem}

\begin{proof}
Assume $\sfname{HypothesesH'}$; in particular its component \ref{pf:vwin-he:H1} holds. We use the following nine
results; for each we say what it states, exactly in the form in which it is used.
\begin{enumerate}[label=(\alph*)]
\item\label{pf:vwin-he:a} \ref{pf:vwin-he:H1} implies that $\sfname{LemmaK}(C_K)$ holds for some $C_K\in\R$
(\lref{UpperTheory.V.ZK.lemmaK_litN_alt}; this is \cref{thm:window:lemmaK}).
\item\label{pf:vwin-he:b} $\sfname{HypothesesH'}$ implies $(\mathrm{U}_{\mathrm{ord}})$ and $(\mathrm{U}_i)$
(\lref{UpperTheory.V.Z.Win.Dn.ends0_alt}).
\item\label{pf:vwin-he:c} $\sfname{HypothesesH'}$ implies $(\mathrm{Sp})$
(\lref{UpperTheory.V.HV.DescentVanishingReductions.sp4In0_alt}).
\item\label{pf:vwin-he:d} $\sfname{HypothesesH'}$ implies $\sfname{KatoTwoW}$
(\lref{UpperTheory.V.Z.Kato.WeilSphericalRecords.katoTwoW_of_printsH}; this is \cref{thm:kato}).
\item\label{pf:vwin-he:e} \ref{pf:vwin-he:H1} and $(\mathrm{U}_{\mathrm{ord}})$ imply $(\mathrm{U}_{\mathrm{ur}})$
(\lref{UpperTheory.V.HV.Core0.automorphyUnramified_of_litN_of_gdEnd}).
\item\label{pf:vwin-he:f} $\sfname{HypothesesH'}$, the existence of $C_K$ with $\sfname{LemmaK}(C_K)$,
$\sfname{KatoTwoW}$ and $(\mathrm{U}_{\mathrm{ur}})$ imply $\sfname{CountI}$
(\lref{UpperTheory.V.Z.J.count_i_of_printed_bound_k_kato_automorphy}; this is \cref{thm:window:countI}).
\item\label{pf:vwin-he:g} $\sfname{HypothesesH'}$, the existence of $C_K$ with $\sfname{LemmaK}(C_K)$, and
$(\mathrm{Sp})$ imply $\sfname{CountII}$ and $\sfname{CountIII}$
(\lref{UpperTheory.V.HV.DescentVanishingReductions.countII_and_countIII_of_printsH_of_sp4In}).
\item\label{pf:vwin-he:h} $\sfname{HypothesesH'}$, the existence of $C_K$ with $\sfname{LemmaK}(C_K)$, and
$(\mathrm{U}_i)$ imply $\sfname{CountV}$ (\lref{UpperTheory.V.Z.J.count_v_of_printed_bound_k_reduced_unitary_lift}).
\item\label{pf:vwin-he:i} \ref{pf:vwin-he:H1}, the existence of $C_K$ with $\sfname{LemmaK}(C_K)$, $\sfname{CountI}$,
$\sfname{CountII}$, $\sfname{CountIII}$ and $\sfname{CountV}$ imply $\sfname{VWin}$
(\lref{UpperTheory.V.HV.Core.of_four_alt}).
\end{enumerate}
By \ref{pf:vwin-he:a} there is $C_K$ with $\sfname{LemmaK}(C_K)$. By \ref{pf:vwin-he:b}, \ref{pf:vwin-he:c} and
\ref{pf:vwin-he:d} the statements $(\mathrm{U}_{\mathrm{ord}})$, $(\mathrm{U}_i)$, $(\mathrm{Sp})$ and
$\sfname{KatoTwoW}$ hold, and then $(\mathrm{U}_{\mathrm{ur}})$ holds by \ref{pf:vwin-he:e}. Now \ref{pf:vwin-he:f}
gives $\sfname{CountI}$, \ref{pf:vwin-he:g} gives $\sfname{CountII}$ and $\sfname{CountIII}$, and \ref{pf:vwin-he:h}
gives $\sfname{CountV}$. Finally \ref{pf:vwin-he:i} gives $\sfname{VWin}$.
\end{proof}

\begin{remark}[On the structure of the reduction]\label{pf:vwin-he:rem-structure}
(1) The package is used in full only through the cited results \ref{pf:vwin-he:b}--\ref{pf:vwin-he:d} and
\ref{pf:vwin-he:f}--\ref{pf:vwin-he:h}; the steps \ref{pf:vwin-he:a}, \ref{pf:vwin-he:e}, \ref{pf:vwin-he:i} use
only its component \ref{pf:vwin-he:H1}.

(2) The statement $\sfname{CountZ}$ of \cref{prop:window:propZ}, the bounds for the $2$-part and the bad layers
(\cref{thm:window:bab}) and the minus part (\cref{prop:psibound}) are not hypotheses of \ref{pf:vwin-he:i}: they are
established in its proof from \ref{pf:vwin-he:H1} and $\sfname{LemmaK}$. Accordingly, the plus-part inputs which
require automorphy are exactly the four counts.

(3) The three- and four-dimensional automorphy statements enter only through
\ref{pf:vwin-he:f}--\ref{pf:vwin-he:h}: $(\mathrm{Sp})$ is used for cases II and III, $(\mathrm{U}_{\mathrm{ur}})$
(deduced from $(\mathrm{U}_{\mathrm{ord}})$) together with $\sfname{KatoTwoW}$ for case I, and $(\mathrm{U}_i)$ for
case V.
\end{remark}

\subsection{Lemma M: cotangent lengths of graded algebras}\label{sec:pf:vwin-he:lemmaM}

In the method of \cref{sec:window:eis}, the length of a Selmer-type group is compared with a graded piece of the
cotangent module of a deformation ring at the Eisenstein point, and the latter is bounded by valuations of a trace
difference at the characteristic-zero points (step (ii) there). The deformation rings carry an action of the
multiplicative group by conjugation with a cocharacter, i.e.\ a $\Z$-grading, the Eisenstein point is fixed by it,
and the other points move in orbits. Lemma M is the abstract form of this step.

\subsubsection{The setting}

Throughout this subsection, $\OO$ is a complete discrete valuation ring of characteristic $0$ with uniformizer
$\varpi$ and fraction field $\mathrm{Frac}(\OO)$, $K$ is an algebraically closed field containing $\OO$ (i.e.\ an
$\OO$-algebra with $\OO\to K$ injective), and $v:K\to\R$ is a \emph{normalized valuation}: $v(ab)=v(a)+v(b)$ and
$v(a+b)\ge\min(v(a),v(b))$ whenever $a$, $b$ and (in the second case) $a+b$ are non-zero, and $v(c)=\ord_\OO(c)$ for
$0\ne c\in\OO$. A local ring $R$ with an $\OO$-algebra structure is \emph{complete over $\OO$} if it is Noetherian,
complete for its maximal ideal, $\OO\to R$ is a local homomorphism, and every element of $R$ is congruent to an
element of $\OO$ modulo the maximal ideal.

\begin{definition}[Weighted rings and their points]\label{pf:vwin-he:wring}
A \emph{weighted ring} is a $\Z$-graded $\OO$-algebra $A=\bigoplus_{j\in\Z}A_j$ (the $A_j$ being $\OO$-submodules
with $A_iA_j\subset A_{i+j}$) together with an $\OO$-algebra homomorphism $x:A\to\OO$ such that: $A_0$ is complete
over $\OO$; $A$ is generated as an $\OO$-algebra by $A_0$ and finitely many homogeneous elements of degrees
$-2,-1,1,2$; and $x(A_j)=0$ for $j\ne0$. A \emph{$K$-point} is an $\OO$-algebra homomorphism $\xi:A\to K$; it lies
on an ideal $I$ if $\xi(I)=0$. The group $K^\times$ acts on $K$-points by $(t\cdot\xi)(f)=t^j\xi(f)$ for $f\in A_j$.
A $K$-point is \emph{negatively trivial} (resp.\ \emph{positively trivial}) if it vanishes on $A_j$ for all $j<0$
(resp.\ $j>0$), and \emph{two-signed} if it is neither. A \emph{point derivation} at $\xi$ is an $\OO$-linear map
$\delta:A\to K$ with $\delta(fg)=\xi(f)\delta(g)+\xi(g)\delta(f)$.
\formal{UpperTheory.V.Alg.IsWRing}
\end{definition}

Two-signed points are never fixed by $K^\times$, and the restriction of a $K$-point to $A_0$ only depends on its
orbit.

\begin{definition}[The saturated cotangent lattice]\label{pf:vwin-he:cot}
Let $(A,x)$ be a weighted ring and $\mathfrak m_x=\ker x$. Let $\vwinhesat(x)=\{f\in A: cf\in\mathfrak m_x^2\
\text{for some}\ 0\ne c\in\OO\}$, an $\OO$-submodule of $A$, and for $j\in\Z$ let
$\vwinhecot_j(x)\subset A/\vwinhesat(x)$ be the image of $A_j\cap\mathfrak m_x$. For a family $F=(F_i)$ of elements
of $A$ let $\langle F\rangle\subset A/\vwinhesat(x)$ be the $\OO$-span of their images. We say that $F$
\emph{spans the cotangent space in degree $j$} if for every $p\in\vwinhecot_j(x)$ there is $0\ne c\in\OO$ with
$cp\in\langle F\rangle$, and we put
\[
 \vwinhecotlen_j(x;F)=\len_\OO\Bigl(\vwinhecot_j(x)\big/\bigl(\vwinhecot_j(x)\cap\langle F\rangle\bigr)\Bigr)\in\N\cup\{\infty\}.
\]
\formallabel{UpperTheory.V.Alg.cotLen}
\end{definition}

Since $\vwinhesat(x)$ is saturated, $A/\vwinhesat(x)$ is a torsion-free $\OO$-module; $\vwinhecot_{-1}(x)$ is the
lattice of degree $-1$ in the cotangent space $\mathfrak m_x/\mathfrak m_x^2$ modulo torsion.

\begin{definition}[Smooth points and charts]\label{pf:vwin-he:chart}
Let $(A,x)$ be a weighted ring, $\tau$ a finite set with a weight function $\mathrm{wt}:\tau\to\Z$, and
$\mathrm{Frac}(\OO)[[\tau]]$ the ring of formal power series in variables indexed by $\tau$; the weight of a monomial
$\prod_iX_i^{\nu_i}$ is $\sum_i\nu_i\mathrm{wt}(i)$. A \emph{chart} at $x$ is an $\OO$-algebra homomorphism
$\phi:A\to\mathrm{Frac}(\OO)[[\tau]]$ such that: $\phi(A_j)$ consists of power series all of whose monomials have
weight $j$; the constant term of $\phi(f)$ is $x(f)$; for each $i\in\tau$ there is $f\in A$ such that the linear
part of $\phi(f)$ is a non-zero multiple of $X_i$; and every $f\in\mathfrak m_x$ such that $\phi(f)$ has no linear
terms lies in $\vwinhesat(x)$. For $n:\Z\to\N$, the point $x$ is \emph{smooth with tangent dimensions $n$} if there
is a chart with $\#\{i:\mathrm{wt}(i)=j\}=n(j)$ for all $j$. A \emph{reduced chart} is a homomorphism $\phi$ with the
first three properties and all of whose variables have weights in $\{-2,-1,1,2\}$.
\formal{UpperTheory.V.Alg.SmoothPt}
\end{definition}

\begin{definition}[Equations and the conditions M1--M5]\label{pf:vwin-he:eqns}
Let $(A,x)$ be a weighted ring and $r,s,r',s',m\in\N$. A \emph{system of equations of type $(r,s,r',s',m)$} consists
of elements $F_0,\dots,F_r\in A_{-1}$, $G_1,\dots,G_s\in A_{-2}$, $E_1,\dots,E_{r'}\in A_1$,
$E'_1,\dots,E'_{s'}\in A_2$ and $Z_1,\dots,Z_m\in A_0\cap\mathfrak m_x$; let $I\subset A$ be the ideal they
generate. Let $n_{r,s,r',s',m}:\Z\to\N$ take the values $r,s,r'+1,s'+1,m$ at $-1,-2,1,2,0$ and $0$ elsewhere. For
$a\in A$ consider the conditions:
\begin{enumerate}[label=(M\arabic*)]
\item $r\ge1$; $x$ is smooth with tangent dimensions $n_{r,s,r',s',m}$; and the $F_i$, the $G_i$ and the $Z_i$
span the cotangent space in degrees $-1$, $-2$ and $0$ respectively;
\item the image of $A_0$ in $A/I$ is a finitely generated $\OO$-module;
\item for every $K$-point $x_1\ne x$ on $I$ which is both negatively and positively trivial, there do not exist both a
$K$-point on $I$ which is negatively but not positively trivial and agrees with $x_1$ on $A_0$, and a $K$-point on
$I$ which is positively but not negatively trivial and agrees with $x_1$ on $A_0$;
\item at every two-signed $K$-point $\xi$ on $I$, any two point derivations at $\xi$ vanishing on $I$ are linearly
dependent over $K$;
\item $a\in A_0$, $x(a)=0$, and $\xi(a)\ne0$ for every two-signed $K$-point $\xi$ on $I$.
\end{enumerate}
A finite set $T$ of two-signed $K$-points on $I$ \emph{covers} $I$ if every two-signed $K$-point on $I$ lies in the
$K^\times$-orbit of an element of $T$.
\formal{UpperTheory.V.Alg.Eqns}
\end{definition}

Condition (M4) says that at a two-signed point the tangent space of $\Spec(A/I)$ over $K$ has dimension at most one,
the dimension of the orbit of the point. Condition (M3) concerns the fixed points other than $x$. Note that there
are $r+1$ equations $F_i$ of degree $-1$, one more than the tangent dimension $r$ of $x$ in that degree.

\begin{theorem}[Lemma M]\label{pf:vwin-he:lemmaM}
Let $(A,x)$ be a weighted ring, $(F,G,E,E',Z)$ a system of equations of type $(r,s,r',s',m)$ generating the ideal
$I$, and $a\in A$, and assume \textup{(M1)}--\textup{(M5)}. Then a finite set covering $I$ exists,
$\vwinhecotlen_{-1}(x;F)<\infty$, and for every finite set $T$ covering $I$,
\[
 \vwinhecotlen_{-1}(x;F)\ \le\ 2\sum_{\xi\in T}v(\xi(a)).
\]
\formallabel{UpperTheory.V.Alg.LemmaM}
\formallabel{UpperTheory.V.Alg.M.main_holds}
\end{theorem}

The theorem holds for every algebraically closed $K\supset\OO$ and every normalized valuation $v$ on $K$. The terms
$v(\xi(a))$ are non-negative (\cref{pf:vwin-he:vnn} below), and the sum does not depend on
the choice of the representatives in the orbits, since $\xi(a)$, $a\in A_0$, is constant on orbits.

\subsubsection{Reduction to a normalized system}

\begin{lemma}[Elementary divisors]\label{pf:vwin-he:sm}
Let $(A,x)$ be a weighted ring, $F_0,\dots,F_r\in A_{-1}$, and assume that $\vwinhecot_{-1}(x)$ is a finitely
generated $\OO$-module of rank $r$ and that the $F_i$ span the cotangent space in degree $-1$. Then there are
$F'_0,\dots,F'_r\in A$ with the same $\OO$-span as $F_0,\dots,F_r$, elements $y_0,\dots,y_{r-1}\in A_{-1}$ and
integers $c_0,\dots,c_{r-1}\ge0$ such that $F'_i-\varpi^{c_i}y_i\in\vwinhesat(x)$ for $i<r$,
$F'_r\in\vwinhesat(x)$, and $\vwinhecotlen_{-1}(x;F)=\sum_{i<r}c_i$.
\formallabel{UpperTheory.V.Alg.M.sm_holds}
\end{lemma}

\begin{proof}
Write $P=\vwinhecot_{-1}(x)$; it is torsion-free, hence free of rank $r$. Since $x(A_{-1})=0$, the images of the
$F_i$ lie in $P$; let $N=P\cap\langle F\rangle$ be their span. As $F$ spans the cotangent space, $P/N$ is torsion,
so $N$ has rank $r$. By the elementary divisor theorem over the discrete valuation ring $\OO$ there are a basis
$b_0,\dots,b_{r-1}$ of $P$ and integers $c_i\ge0$ such that $\varpi^{c_0}b_0,\dots,\varpi^{c_{r-1}}b_{r-1}$ is a basis
of $N$; then $P/N\cong\bigoplus_i\OO/\varpi^{c_i}$ has length $\sum_ic_i$. Let $\Psi:\OO^{r+1}\to P$ be the map
$e_i\mapsto$ image of $F_i$; its image is $N$, of rank $r$, so its kernel has rank one. Choose $w_i\in\OO^{r+1}$ with
$\Psi(w_i)=\varpi^{c_i}b_i$ ($i<r$) and a generator $w_r$ of $\ker\Psi$. Then $w_0,\dots,w_r$ generate $\OO^{r+1}$
(for $u\in\OO^{r+1}$ write $\Psi(u)=\sum a_i\varpi^{c_i}b_i$; then $u-\sum a_iw_i\in\ker\Psi$), hence form a basis.
Put $F'_i=\sum_kw_{i,k}F_k$; these have the same span as the $F_k$. Let $y_i\in A_{-1}\cap\mathfrak m_x$ be lifts of
$b_i$. Then the image of $F'_i$ in $A/\vwinhesat(x)$ is $\varpi^{c_i}b_i$ for $i<r$ and $0$ for $i=r$, which are the
remaining assertions.
\end{proof}

\begin{lemma}[Homogeneity]\label{pf:vwin-he:hom}
For a system of equations as in \cref{pf:vwin-he:eqns}, the ideal $I$ and the ideal $I_2$ generated by the
$F_i^2$, $G_i$, $E_i^2$, $E'_i$ and $Z_i$ are homogeneous, and $I_2\subset I$.
\formallabel{UpperTheory.V.Alg.M.hom_holds}
\end{lemma}

\begin{proof}
Both ideals are generated by homogeneous elements, and $F_i^2=F_i\cdot F_i$, $E_i^2=E_i\cdot E_i\in I$.
\end{proof}

The following four lemmas normalize the data; their proofs are given in \cref{sec:pf:vwin-he:pf-chart} and
\cref{sec:pf:vwin-he:pf-points}. The notion of a finite set covering an ideal (\cref{pf:vwin-he:eqns}) makes sense for
every ideal of $A$.

\begin{lemma}[Rank of the cotangent lattice]\label{pf:vwin-he:rk}
Let $(A,x)$ be a weighted ring which is smooth at $x$ with tangent dimensions $n$. Then $\vwinhecot_{-1}(x)$ is a
finitely generated $\OO$-module of rank $n(-1)$.
\formallabel{UpperTheory.V.Alg.M.rk_holds}
\end{lemma}

\begin{lemma}[Elimination of the variables of weight zero]\label{pf:vwin-he:cut}
Let $(A,x)$ be a weighted ring and $(F,G,E,E',Z)$ a system of equations of type $(r,s,r',s',m)$ satisfying
\textup{(M1)}. Then there are a finite set $\tau$ with a weight function and a reduced chart
$\phi:A\to\mathrm{Frac}(\OO)[[\tau]]$ such that $\phi(Z_i)=0$ for all $i$ and
$\#\{i:\mathrm{wt}(i)=j\}=n_{r,s,r',s',0}(j)$ for all $j\in\Z$.
\formallabel{UpperTheory.V.Alg.M.cut_holds}
\end{lemma}

\begin{lemma}[Non-vanishing on the chart]\label{pf:vwin-he:nz}
Let $(A,x)$ be a weighted ring, $(F,G,E,E',Z)$ a system of equations with ideal $I$ such that the $F_i$ span the
cotangent space in degree $-1$, and $\phi:A\to\mathrm{Frac}(\OO)[[\tau]]$ a reduced chart with at least one variable
of weight $-1$ and at least one variable of weight $1$. Then for every $a\in A_0$ there is $a'\in A_0$ with
$a'-a\in I$, $x(a')=x(a)$ and $\phi(a')\ne0$.
\formallabel{UpperTheory.V.Alg.M.nz_holds}
\end{lemma}

\begin{lemma}[Existence of covering sets]\label{pf:vwin-he:fin}
Let $(A,x)$ be a weighted ring and $J\subset A$ a homogeneous ideal such that the image of $A_0$ in $A/J$ is a
finitely generated $\OO$-module and \textup{(M4)} holds with $J$ in place of $I$. Then a finite set covering $J$
exists.
\formallabel{UpperTheory.V.Alg.M.fin_holds}
\end{lemma}

\begin{theorem}[Core inequality]\label{pf:vwin-he:core}
Let $(A,x)$ be a weighted ring, $(F,G,E,E',Z)$ a system of equations of type $(r,s,r',s',m)$ with ideal $I$, and
$\phi:A\to\mathrm{Frac}(\OO)[[\tau]]$ a reduced chart with $\phi(Z_i)=0$ for all $i$ and
$\#\{i:\mathrm{wt}(i)=j\}=n_{r,s,r',s',0}(j)$ for all $j$. Assume $r\ge1$, and let $y_0,\dots,y_{r-1}\in A_{-1}$ and
$c_0,\dots,c_{r-1}\in\N$ satisfy $F_i-\varpi^{c_i}y_i\in\vwinhesat(x)$ for $i<r$ and $F_r\in\vwinhesat(x)$. Assume
that the $F_i$ and the $G_i$ span the cotangent space in degrees $-1$ and $-2$, that \textup{(M2)}, \textup{(M3)},
\textup{(M4)} hold, and that $a$ satisfies \textup{(M5)} and $\phi(a)\ne0$. Then for every finite set $T$ covering $I$,
\[
 \sum_{i<r}c_i\ \le\ 2\sum_{\xi\in T}v(\xi(a)).
\]
\formallabel{UpperTheory.V.Alg.M.core_holds}
\end{theorem}

\begin{proof}[Proof of \cref{pf:vwin-he:lemmaM}, assuming \cref{pf:vwin-he:core}]
By (M1) and \cref{pf:vwin-he:rk}, $\vwinhecot_{-1}(x)$ is finitely generated of rank
$r$, and $x(A_{-1})=0$. \Cref{pf:vwin-he:sm} gives $F'$, $y$ and $c$ with
$\vwinhecotlen_{-1}(x;F)=\sum_ic_i<\infty$. Since the $F'_i$ have the same $\OO$-span as the $F_i$, they generate the
same ideal and $\langle F'\rangle=\langle F\rangle$; so replacing $F$ by $F'$ does not change $I$, the conditions
(M1)--(M5), or $\vwinhecotlen_{-1}$. We may therefore assume $F=F'$. By \cref{pf:vwin-he:cut} there is a reduced
chart $\phi$ with $\phi(Z_i)=0$ and $r\ge1$ variables of weight $-1$ and $r'+1\ge1$ of weight $1$. By
\cref{pf:vwin-he:nz}, applied to $a\in A_0$, there is $a'\in A_0$ with $a'-a\in I$, $x(a')=x(a)=0$ and
$\phi(a')\ne0$. Every $K$-point on $I$ takes the same value at $a'$ and at $a$, so $a'$ satisfies (M5) and
$\sum_{\xi\in T}v(\xi(a'))=\sum_{\xi\in T}v(\xi(a))$ for every $T$. A finite set covering $I$ exists by
\cref{pf:vwin-he:hom} and \cref{pf:vwin-he:fin}, and the inequality follows from \cref{pf:vwin-he:core} applied to
$a'$.
\end{proof}

\subsubsection{Euler characteristics on multigraded algebras}

We recall the calculus of Euler characteristics of Koszul type used in the proof of \cref{pf:vwin-he:core}. Let
$k\ge1$; the group $\Z^k$ is partially ordered componentwise, $e_0,\dots,e_{k-1}$ are its unit vectors, and a
statement about $d\in\Z^k$ holds \emph{for large $d$} if there is $d_0$ such that it holds for all $d\ge d_0$.

\begin{definition}[Standard settings, graded modules and Euler characteristics]\label{pf:vwin-he:std}
A \emph{standard setting} is a Noetherian local ring $B$ and a $\Z^k$-graded $B$-algebra $R=\bigoplus_dR_d$ (the
$R_d$ being $B$-submodules) such that $R_d=0$ unless $d\ge0$, $R_0$ is the image of $B$, each $R_{e_i}$ is a
finitely generated $B$-module, and $R_{d+e_i}=R_dR_{e_i}$ for $d\ge0$. Its \emph{irrelevant ideal} is
$\vwinheIrr=\prod_iR\cdot R_{e_i}$. A \emph{graded module} is a finitely generated $R$-module $M$ with $B$-submodules
$M_d$ such that $M=\bigoplus_dM_d$ and $R_dM_{d'}\subset M_{d+d'}$; for $s\in R$ we write $M/sM$ for the quotient
with the induced grading, and $M[s](\delta)$ for the submodule $\{y:sy=0\}$ with grading $(M[s](\delta))_d=M[s]\cap
M_{d-\delta}$. A \emph{valid list} is a finite list $ts=((\delta_1,s_1),\dots,(\delta_n,s_n))$ with $\delta_i\ge0$
and $s_i\in R_{\delta_i}$; $\vwinhesecs(ts)$ is the ideal generated by the $s_i$. The \emph{Euler characteristic}
$\chi(ts,M)(d)\in\Z$ is defined recursively by $\chi((),M)(d)=\len_BM_d$ (with the convention that an infinite
length counts as $0$) and
\[
 \chi\bigl(((\delta,s))\cdot ts,M\bigr)(d)=\chi(ts,M/sM)(d)-\chi(ts,M[s](\delta))(d).
\]
A list $ts$ is \emph{admitted} for an ideal $J\subset R$ if $(\mathfrak m_BR\cdot\vwinheIrr)^n\subset
J+\vwinhesecs(ts)$ for some $n$.
\formal{Formalism.MainNotions.V.Pf.V2.Calc.chi}
\end{definition}

For $ts=((\delta_1,s_1),\dots,(\delta_n,s_n))$ and a module $M$ killed by all $s_i$, unwinding the definition gives the
alternating sum $\sum_{S\subset\{1..n\}}(-1)^{|S|}\len M_{d-\sum_{i\in S}\delta_i}$; in general $\chi(ts,M)(d)$ is an
Euler characteristic of Koszul type, built from the kernels and cokernels of the multiplications by the $s_i$. We use
the following rules, for a standard setting, graded modules $M$, valid lists and large $d$:
\begin{enumerate}[label=(E\arabic*)]
\item\label{pf:vwin-he:E-gr} submodules generated by homogeneous elements, quotients by them, shifts, $M/sM$ and
$M[s](\delta)$ for $s\in R_\delta$ are graded modules; the pieces $M_d$ are finitely generated $B$-modules; and if
$()$ is admitted for the annihilator $\vwinheAnn(M)$, then $\len_BM_d<\infty$ for large $d$
(\lref{Formalism.MainNotions.V.Pf.V2.Calc.ruleGr} and \lref{Formalism.MainNotions.V.Pf.V2.Calc.ruleFin});
\item\label{pf:vwin-he:E-iso} $\chi(ts,M)=\chi(ts,M')$ if $M\cong M'$ as graded modules
(\lref{Formalism.MainNotions.V.Pf.V2.Calc.ruleIso});
\item\label{pf:vwin-he:E-nzd} if $s$ is a non-zero-divisor on $M$, then $\chi(((\delta,s))\cdot ts,M)=\chi(ts,M/sM)$
(\lref{Formalism.MainNotions.V.Pf.V2.Calc.ruleNzd});
\item\label{pf:vwin-he:E-perm} if $ts$ is admitted for $\vwinheAnn(M)$ and $ts'$ is a permutation of $ts$, then
$\chi(ts,M)(d)=\chi(ts',M)(d)$ for large $d$
(\lref{Formalism.MainNotions.V.Pf.V2.Calc.euler_rule_two_of_graded_finite_iso}, together with the additivity of
$\chi$ in short exact sequences, \lref{Formalism.MainNotions.V.Pf.V2.Calc.euler_additivity_of_rules});
\item\label{pf:vwin-he:E-ub} if $ts\cdot us$ is admitted for $\vwinheAnn(M)$, then $\chi(ts\cdot us,M)(d)\le
\chi(us,M/\vwinhesecs(ts)M)(d)$ for large $d$ (\lref{Formalism.MainNotions.V.Pf.V2.Calc.ruleUB});
\item\label{pf:vwin-he:E-mono} if $0\to M'\to M\to M''\to0$ is an exact sequence of graded modules and $ts$ is admitted
for $\vwinheAnn(M)$, then $\chi(ts,M'')(d)\le\chi(ts,M)(d)$ for large $d$
(\lref{Formalism.MainNotions.V.Pf.V2.Calc.ruleMono});
\item\label{pf:vwin-he:E-quot} for an ideal $J\subset R$, the successive quotient of $R/J$ by the elements of $ts$ is
isomorphic, as a graded module, to $R/(J+\vwinhesecs(ts))$
(\lref{UpperTheory.V.Alg.M.qi_holds});
\item\label{pf:vwin-he:E-cong} if $ts$ and $ts'$ have the same degrees and $s_i-s'_i\in\vwinheAnn(M)$ for all $i$,
then $\chi(ts,M)=\chi(ts',M)$;
\item\label{pf:vwin-he:E-single} for $\beta\in B$, $\chi(((0,\beta)),M)(d)=\theta(\beta,M_d)$, where for a $B$-module
$N$ we put $\theta(\beta,N)=\len_B(N/\beta N)-\len_BN[\beta]$ (finite lengths understood as above);
\item\label{pf:vwin-he:E-prod} if $s\in R_0$, $s'\in R_\delta$, the list $((0,s),(\delta,s'))\cdot ts$ is valid and
$((\delta,ss'))\cdot ts$ is admitted for $\vwinheAnn(M)$, then for large $d$
\[
 \chi\bigl(((\delta,ss'))\cdot ts,M\bigr)(d)=\chi\bigl(((0,s))\cdot ts,M\bigr)(d)+\chi\bigl(((\delta,s'))\cdot ts,M\bigr)(d)
\]
(\lref{Formalism.MainNotions.V.Pf.V2.Calc.euler_rule_product_zero});
\item\label{pf:vwin-he:E-small} if $()$ is admitted for $\vwinheAnn(M)$ and $ts=((\delta_1,s_1),\dots,(\delta_n,s_n))$
is valid, then $\chi(ts,M)(d)=\sum_{S\subset\{1,\dots,n\}}(-1)^{|S|}\len_BM_{d-\sum_{i\in S}\delta_i}$ for large $d$
(\lref{Formalism.MainNotions.V.Pf.V2.Calc.ruleSmall});
\item\label{pf:vwin-he:E-pos} if $ts$ is admitted for $\vwinheAnn(M)$, then $\chi(ts,M)(d)\ge0$ for large $d$
(\lref{Formalism.MainNotions.V.Pf.V2.Calc.rulePos}).
\end{enumerate}
\begin{lemma}\label{pf:vwin-he:E89}
Rules \ref{pf:vwin-he:E-cong} and \ref{pf:vwin-he:E-single} hold, for every graded module $M$ (and every $d$).
\formallabel{Formalism.MainNotions.V.Pf.V2.Calc.Cong.chi_cong}
\formallabel{Formalism.MainNotions.V.Pf.V2.Calc.chi_single0}
\end{lemma}

\begin{proof}
For \ref{pf:vwin-he:E-cong}, induct on the length of the lists: if $s-s'$ kills $M$, then $sM=s'M$ and $M[s]=M[s']$, so
$M/sM=M/s'M$ and $M[s](\delta)=M[s'](\delta)$, and the annihilators of these modules contain $\vwinheAnn(M)$, so the
induction hypothesis applies to the remaining entries. For \ref{pf:vwin-he:E-single}, the definition gives
$\chi(((0,\beta)),M)(d)=\len_B(M/\beta M)_d-\len_B(M[\beta](0))_d$, and since $\beta$ has degree $0$ and $M$ is the
direct sum of its pieces, $(M/\beta M)_d=M_d/\beta M_d$ and $(M[\beta](0))_d=M_d[\beta]$.
\end{proof}

\begin{definition}[The Rees algebra of a system of equations]\label{pf:vwin-he:rees}
Let $(A,x)$ be a weighted ring and $(F,G,E,E',Z)$ a system of equations. Put $B=A_0$ and
let $N_0=A_{-2}$, $N_1=A_2$ and $N_2=B\,F_r^2+A_{-2}A_{-2}A_2$, which are $B$-submodules of $A$, and for
$\mu\in\N^3$ let $N^\mu=N_0^{\mu_0}N_1^{\mu_1}N_2^{\mu_2}$ (products of $B$-submodules of $A$, with $N_i^0=B$). The
\emph{Rees algebra} is
\[
 R=\bigoplus_{\mu\in\N^3}N^\mu\,t^\mu\ \subset\ A[t_0,t_1,t_2],
\]
graded by $\mu\in\N^3\subset\Z^3$; for $p\in R$ we write $p_\mu\in N^\mu$ for the coefficient of $t^\mu$. For
$f\in N_i$ write $f t_i\in R_{e_i}$. For an ideal $J\subset A$ let
$J^R\subset R$ be the ideal of elements all of whose coefficients lie in $J$. The list of sections is
\[
 ts=\bigl((e_0,F_i^2t_0)_{i<r},\ (e_0,G_it_0)_{i\le s},\ (e_1,E_i^2t_1)_{i\le r'},\ (e_1,E'_it_1)_{i\le s'},\
 (e_2,F_r^2t_2)\bigr),
\]
and $ts_{\mathrm F}$, $ts_{\mathrm R}$ denote its first part (the $r$ entries $F_i^2t_0$) and the rest.
\formal{UpperTheory.V.Alg.M.mR}
\end{definition}

The entries of $ts$ lie in the right degrees since $F_i^2\in A_{-2}$ and $E_i^2\in A_2$.

\begin{lemma}[Graded pieces of weighted rings]\label{pf:vwin-he:fg}
Let $(A,x)$ be a weighted ring. Then every $A_j$ is a finitely generated $A_0$-module, and for $n\in\N$ we have
$A_{2n}=(A_2)^n$ and $A_{-2n}=(A_{-2})^n$, as products of $A_0$-submodules of $A$ (with $(A_{\pm2})^0=A_0$).
\formallabel{UpperTheory.V.Alg.M.fg_holds}
\formallabel{UpperTheory.V.Alg.M.oz_two_mul_eq_pow}
\end{lemma}

\begin{proof}
The ring $A_0$ is Noetherian, and $A$ is generated as an $A_0$-algebra by finitely many homogeneous elements, so $A$
is Noetherian by Hilbert's basis theorem. Hence the ideal $A\cdot A_j$ is generated by finitely many elements
$t_1,\dots,t_k\in A_j$. For $y\in A_j$ write $y=\sum_if_it_i$ with $f_i\in A$; taking components of degree $j$ gives
$y=\sum_i(f_i)_0t_i$, where $(f_i)_0\in A_0$ is the component of degree $0$ of $f_i$. So $t_1,\dots,t_k$ generate
$A_j$ over $A_0$.

The inclusions $(A_{\pm2})^n\subset A_{\pm2n}$ hold as $A_iA_j\subset A_{i+j}$. Since $A$ is generated as an
$\OO$-algebra by $A_0$ and homogeneous elements of degrees $\pm1,\pm2$, each $A_j$ is spanned over $\OO$ by products
$z_1\cdots z_k$ of homogeneous elements $z_i$ of degrees $w_i\in\{-2,\dots,2\}$ with $\sum_iw_i=j$. If $j\ge2$, such
a product has a factor of degree $2$ or two factors of degree $1$, since otherwise $\sum_iw_i\le1$; collecting these
factors writes it as $uu'$ with $u\in A_2$ and $u'\in A_{j-2}$. Hence $A_j\subset A_2A_{j-2}$ for $j\ge2$, and
$A_{2n}\subset(A_2)^n$ follows by induction on $n$. The case of $A_{-2n}$ is symmetric.
\end{proof}

\begin{lemma}[Rees algebras are standard]\label{pf:vwin-he:stdR}
Let $B$ be a Noetherian local ring, $A$ a $B$-algebra, $k\ge1$ and $N_0,\dots,N_{k-1}$ finitely generated
$B$-submodules of $A$. Then $B$ and the Rees algebra $R=\bigoplus_{\mu\in\N^k}N^\mu t^\mu\subset
A[t_0,\dots,t_{k-1}]$, where $N^\mu=\prod_iN_i^{\mu_i}$, graded by $\Z^k$ with $R_d=0$ for $d\notin\N^k$, form a
standard setting.
\formallabel{UpperTheory.V.Alg.M.std_holds}
\end{lemma}

\begin{proof}
Since $N^\mu N^{\mu'}=N^{\mu+\mu'}$, $R$ is a $\Z^k$-graded $B$-subalgebra of $A[t_0,\dots,t_{k-1}]$ whose elements
of degree $\mu$ are exactly the $ft^\mu$ with $f\in N^\mu$. Thus $R_d=0$ unless $d\ge0$; $R_0=B\cdot1$ is the image of
$B$; $R_{e_i}=N_it_i\cong N_i$ is finitely generated; and for $\mu\ge0$ we have
$R_{\mu+e_i}=N^{\mu}N_i\,t^{\mu+e_i}=R_\mu R_{e_i}$, since the products $ft^\mu\cdot gt_i$ with $f\in N^\mu$ and
$g\in N_i$ span $N^\mu N_i\,t^{\mu+e_i}$.
\end{proof}

By \cref{pf:vwin-he:fg}, the modules $N_0$, $N_1$ and $N_2$ of \cref{pf:vwin-he:rees} are finitely generated over the
Noetherian local ring $B=A_0$, so $(B,R)$ is a standard setting by \cref{pf:vwin-he:stdR}.

\subsubsection{The two estimates}

For an $\OO$-algebra homomorphism $\sigma:B\to K$ and a $B$-module $M$ let $\vwinheloclen_\sigma(M)$ be the length of
the localization of $M$ at the prime $\ker\sigma$, as a module over $B_{\ker\sigma}$. For an ideal $J\subset A$ and
$w\in\Z$ let $(A/J)_w$ be the image of $A_w$ in $A/J$, a $B$-module; similarly, for an ideal $J'$ of the Rees algebra
$R$ and $d\in\Z^3$, $(R/J')_d$ is the image of $R_d$ in $R/J'$. Let $I_D\subset B=A_0$ be the ideal generated by
the products $fg$ with $f\in A_{-2}$, $g\in A_2$. Recall the ideal $I_2\subset I$ of \cref{pf:vwin-he:hom}.

\begin{lemma}\label{pf:vwin-he:sign}
Let $(A,x)$ be a weighted ring and $\xi:A\to\kappa$ an $\OO$-algebra homomorphism to a field. Then $\xi$ is not
negatively trivial if and only if $\xi(f)\ne0$ for some $f\in A_{-2}$, and not positively trivial if and only if
$\xi(g)\ne0$ for some $g\in A_2$. Consequently a $K$-point $\xi$ is two-signed if and only if $\xi|_{A_0}$ does not
vanish on $I_D$.
\formallabel{UpperTheory.V.Alg.M.sign_holds}
\end{lemma}

\begin{proof}
If $\xi$ vanishes on $A_{-2}$, then $\xi(f)^2=\xi(f^2)=0$ for $f\in A_{-1}$, so $\xi$ vanishes on the homogeneous
generators of negative degree. Every element of $A_j$, $j<0$, is a sum of products of elements of $A_0$ and of the
homogeneous generators of total degree $j$, and each such product has a factor of negative degree; hence $\xi$
vanishes on $A_j$. The converse is clear, and the positive case is symmetric. For the last assertion, $\xi(I_D)\ne0$
if and only if $\xi(f)\xi(g)=\xi(fg)\ne0$ for some $f\in A_{-2}$, $g\in A_2$.
\end{proof}

The next five lemmas are proved in \cref{sec:pf:vwin-he:pf-val,sec:pf:vwin-he:pf-points,sec:pf:vwin-he:pf-rees}.

\begin{lemma}[Non-negativity]\label{pf:vwin-he:vnn}
Let $A=\bigoplus_jA_j$ be a $\Z$-graded $\OO$-algebra with $A_0$ complete over $\OO$, $I\subset A$ an ideal such
that the image of $A_0$ in $A/I$ is a finitely generated $\OO$-module, $\xi:A\to K$ an $\OO$-algebra homomorphism
vanishing on $I$, and $b\in A_0$ with $\xi(b)\ne0$. Then $v(\xi(b))\ge0$.
\formallabel{UpperTheory.V.Alg.M.valuation_nonneg_apply_degreeZero_of_onS}
\end{lemma}

\begin{lemma}[A bound of Herbrand type]\label{pf:vwin-he:th}
Let $B'$ be complete over $\OO$, $M$ a finitely generated $B'$-module such that $B'/\vwinheAnn(M)$ is a finitely
generated $\OO$-module, $a\in B'$, $c\in\N$, and $S$ a finite set of $\OO$-algebra homomorphisms $B'\to K$ such
that $\sigma(a)\ne0$, $v(\sigma(a))\ge0$ and $\vwinheloclen_\sigma(M)\le c$ for $\sigma\in S$, and
$\vwinheloclen_\sigma(M)=0$ for all other $\OO$-algebra homomorphisms $\sigma:B'\to K$. Then, with $\theta$ as in
\ref{pf:vwin-he:E-single} (for $B'$ in place of $B$),
\[
 \theta(a,M)\ \le\ c\sum_{\sigma\in S}v(\sigma(a)).
\]
\formallabel{UpperTheory.V.Alg.M.th_holds}
\end{lemma}

The proof shows that $\theta(a,M)=\sum_{\sigma\in S}\vwinheloclen_\sigma(M)\,v(\sigma(a))$; for a domain $M=D$ which
is finite and torsion-free over $\OO$ this is the classical formula $\len_\OO(D/aD)=\sum_\sigma v(\sigma(a))$, the
sum being over the $\OO$-algebra homomorphisms $\sigma:D\to K$ (\cref{pf:vwin-he:lqs}).

\begin{lemma}[Vanishing of localized lengths]\label{pf:vwin-he:ann}
Let $(A,x)$ be a weighted ring and $(F,G,E,E',Z)$ a system of equations with ideal $I$ such that the $F_i$ and the
$G_i$ span the cotangent space in degrees $-1$ and $-2$ and \textup{(M2)}, \textup{(M3)} hold. Let $R$ and $ts$ be as
in \cref{pf:vwin-he:rees}, and $J'\subset R$ an ideal generated by homogeneous elements which contains
$\vwinhesecs(ts)$ and the elements $Z_i\in B=R_0$. Then for large $d$: $B/\vwinheAnn_B((R/J')_d)$ is a finitely
generated $\OO$-module, and $\vwinheloclen_\sigma((R/J')_d)=0$ for every $\OO$-algebra homomorphism $\sigma:B\to K$
which vanishes on $I_D$ or does not vanish on $B\cap I$.
\formallabel{UpperTheory.V.Alg.M.eventually_mQ_gr_finite_annihilator}
\end{lemma}

\begin{lemma}[Comparison with the weight pieces]\label{pf:vwin-he:loc1}
Let $(A,x)$ be a weighted ring, $(F,G,E,E',Z)$ a system of equations, $R$ and $ts$ as in \cref{pf:vwin-he:rees}, and
$J'\subset R$ an ideal containing $\vwinhesecs(ts)$ and the $Z_i$. Let $\sigma:B\to K$ be an $\OO$-algebra
homomorphism which does not vanish on $I_D$, and $d\in\Z^3$ with $d\ge e_0+e_1+e_2$. Then
\[
 \vwinheloclen_\sigma\bigl((R/J')_d\bigr)\ \le\ \vwinheloclen_\sigma\bigl((A/I_2)_{2(d_1-d_0-d_2)}\bigr).
\]
\formallabel{UpperTheory.V.Alg.M.locLen_mQ_gr_le_locLen_aPiece}
\end{lemma}

\begin{lemma}[Localized lengths at two-signed points]\label{pf:vwin-he:loc3}
Let $(A,x)$ be a weighted ring and $J\subset A$ a homogeneous ideal such that the image of $A_0$ in $A/J$ is a
finitely generated $\OO$-module and \textup{(M4)} holds with $J$ in place of $I$. Let $\sigma:B\to K$ be an
$\OO$-algebra homomorphism which does not vanish on $I_D$. Then $\vwinheloclen_\sigma((A/J)_w)\le1$ for all $w\in\Z$.
\formallabel{UpperTheory.V.Alg.M.locLen_quotientWeightPiece_le_one}
\end{lemma}

We shall also use the following three results of the formal development, whose proofs we do not reproduce.

\begin{lemma}\label{pf:vwin-he:aux}
Let $(A,x)$ be a weighted ring, $(F,G,E,E',Z)$ a system of equations of type $(r,s,r',s',m)$ with ideal $I$, and $R$,
$ts$ as in \cref{pf:vwin-he:rees}.
\begin{enumerate}[label=(\alph*)]
\item\label{pf:vwin-he:aux-loc2} If $\sigma:B\to K$ satisfies $\vwinheloclen_\sigma((A/I)_w)\le1$ for all $w$, then
$\vwinheloclen_\sigma((A/I_2)_w)\le2^{r+1+r'}$ for all $w$
(\lref{UpperTheory.V.Alg.M.locLen_aPiece_ideal2_le_two_pow}).
\item\label{pf:vwin-he:aux-ext} If $J\subset A$ is a homogeneous ideal and $\sigma:B\to K$ vanishes on $B\cap J$,
there is a $K$-point $\xi$ on $J$ with $\xi|_B=\sigma$ (\lref{UpperTheory.V.Alg.M.tsp_holds}).
\item\label{pf:vwin-he:aux-adm} If the $F_i$ and $G_i$ span the cotangent space in degrees $-1$ and $-2$,
\textup{(M2)}, \textup{(M3)} hold and $a$ satisfies \textup{(M5)}, then for every ideal $J\subset R$ containing the
$Z_i$, the list $ts\cdot((0,a))$ is admitted for $J$ (\lref{UpperTheory.V.Alg.M.admissibility_holds}).
\end{enumerate}
\end{lemma}

\begin{proposition}[Upper bound]\label{pf:vwin-he:up}
Let $(A,x)$ be a weighted ring and $(F,G,E,E',Z)$ a system of equations with ideal $I$ such that the $F_i$ and $G_i$
span the cotangent space in degrees $-1$ and $-2$ and \textup{(M2)}--\textup{(M4)} hold; let $a$ satisfy
\textup{(M5)}, $T$ a finite set covering $I$, and $J'\subset R$ an ideal generated by homogeneous elements which
contains $\vwinhesecs(ts)$ and the $Z_i$. Then for large $d$,
\[
 \chi\bigl(((0,a)),R/J'\bigr)(d)\ \le\ 2^{r+1+r'}\sum_{\xi\in T}v(\xi(a)).
\]
\formallabel{UpperTheory.V.Alg.M.up_holds}
\end{proposition}

\begin{proof}
By \cref{pf:vwin-he:sign}, a $K$-point is two-signed if and only if its restriction to $B$ does not vanish on $I_D$.
For $\xi\in T$ we have $\xi(a)\ne0$ by (M5) and $v(\xi(a))\ge0$ by \cref{pf:vwin-he:vnn}. Let
$S=\{\xi|_B:\xi\in T\}$; since the terms are non-negative,
$\sum_{\sigma\in S}v(\sigma(a))\le\sum_{\xi\in T}v(\xi(a))$. By \ref{pf:vwin-he:E-gr} and \cref{pf:vwin-he:ann}, for
large $d$ (we also take $d\ge e_0+e_1+e_2$) the $B$-module
$M_d=(R/J')_d$ is finitely generated, $B/\vwinheAnn(M_d)$ is a finitely generated $\OO$-module, and
$\vwinheloclen_\sigma(M_d)=0$ whenever $\sigma$ vanishes on $I_D$ or not on $B\cap I$. Fix such a $d$; we apply
\cref{pf:vwin-he:th} to $B'=B$, $M=M_d$, $a$, $S$ and $c=2^{r+1+r'}$.

Let $\sigma=\xi|_B\in S$. Then $\sigma(a)\ne0$ and $v(\sigma(a))\ge0$; as $\xi$ is two-signed, $\sigma$ does not vanish
on $I_D$, so by \cref{pf:vwin-he:loc1}, \cref{pf:vwin-he:aux}\ref{pf:vwin-he:aux-loc2} and \cref{pf:vwin-he:loc3}
(with $J=I$, homogeneous by \cref{pf:vwin-he:hom})
\[
 \vwinheloclen_\sigma(M_d)\le\vwinheloclen_\sigma\bigl((A/I_2)_{2(d_1-d_0-d_2)}\bigr)\le2^{r+1+r'}.
\]
Let $\sigma\notin S$ and suppose $\vwinheloclen_\sigma(M_d)\ne0$. Then $\sigma$ does not vanish on $I_D$ and vanishes on
$B\cap I$. By \cref{pf:vwin-he:aux}\ref{pf:vwin-he:aux-ext} there is a $K$-point $\xi$ on $I$ with $\xi|_B=\sigma$;
it is two-signed, so $\xi=t\cdot\xi'$ for some $\xi'\in T$ and $t\in K^\times$. Since $K^\times$ acts trivially on
$A_0=B$, $\sigma=\xi'|_B\in S$, a contradiction. Hence \cref{pf:vwin-he:th} gives
$\theta(a,M_d)\le2^{r+1+r'}\sum_{\sigma\in S}v(\sigma(a))$, and $\chi(((0,a)),R/J')(d)=\theta(a,M_d)$ by
\ref{pf:vwin-he:E-single}.
\end{proof}

For the lower bound we use a filtration on the target of the chart. Let $\phi:A\to\mathrm{Frac}(\OO)[[\tau]]$ be a
reduced chart. The \emph{negative weight} of a monomial $X^\nu=\prod_iX_i^{\nu_i}$ is
$|\nu|_-=\sum_i\nu_i\max(-\mathrm{wt}(i),0)$, and its \emph{positive weight} is
$|\nu|_+=\sum_i\nu_i\max(\mathrm{wt}(i),0)$. Let $\mathrm{Fil}^n$ be the set of power series without monomials of
negative weight $<n$, and let $\mathrm{cnt}(a,b)$ be the number of monomials $\nu$ with $|\nu|_-=a$ and $|\nu|_+=b$;
it is finite since no variable has weight $0$. For a list $L=(\delta_1,\dots,\delta_n)$ in $\Z^k$ and a function $h$
on $\Z^k$ let $\Delta_Lh$ be the iterated difference: $\Delta_{()}h=h$ and
$\Delta_{(\delta)\cdot L}h(d)=\Delta_Lh(d)-\Delta_Lh(d-\delta)$. The following four lemmas are proved in
\cref{sec:pf:vwin-he:pf-chart,sec:pf:vwin-he:pf-count}.

\begin{lemma}[The leading ideal]\label{pf:vwin-he:lk}
Let $(A,x)$ be a weighted ring, $(F,G,E,E',Z)$ a system of equations of type $(r,s,r',s',m)$ with
$F_r\in\vwinhesat(x)$, $\phi:A\to\mathrm{Frac}(\OO)[[\tau]]$ a reduced chart, and $R$ the Rees algebra of
\cref{pf:vwin-he:rees}. Let $\mathfrak q$ be the set of $p\in R$ with $\phi(p_\mu)\in\mathrm{Fil}^{2\mu_0+4\mu_2+1}$
for all $\mu\in\N^3$. Then $\mathfrak q$ is an ideal of $R$ generated by homogeneous elements;
$(\ker\phi)^R\subset\mathfrak q$; every $b\in B$ with $x(b)=0$ lies in $\mathfrak q$; if $0\ne c\in\OO$ and $p\in R$
satisfy $cp\in\mathfrak q$, then $p\in\mathfrak q$; and $f^2t_0-g^2t_0\in\mathfrak q$ for all $f,g\in A_{-1}$ with
$f-g\in\vwinhesat(x)$.
\formallabel{UpperTheory.V.Alg.M.lk_holds}
\end{lemma}

\begin{lemma}[Hilbert function of the leading quotient]\label{pf:vwin-he:hf}
In the situation of \cref{pf:vwin-he:lk}, let $\varpi$ be a uniformizer of $\OO$. Then for all $d\in\N^3$,
\[
 \len_B\bigl((R/\mathfrak q)/\varpi(R/\mathfrak q)\bigr)_d=\mathrm{cnt}(2d_0+4d_2,\,2d_1+2d_2).
\]
\formallabel{UpperTheory.V.Alg.M.hf_holds}
\end{lemma}

\begin{lemma}[A leading coefficient]\label{pf:vwin-he:num}
Let $\tau$ be a finite set with a weight function $\mathrm{wt}$ such that $\#\{i:\mathrm{wt}(i)=j\}=n_{r,s,r',s',0}(j)$
for all $j\in\Z$, where $r\ge1$, and let $\mathrm{cnt}$ be formed with these weights. Let
$L=\bigl((e_0)^{r-1+s},(e_1)^{r'+s'},e_2\bigr)$, exponents denoting repetition. Then for large $d\in\Z^3$,
\[
 \Delta_L\bigl(d\mapsto\mathrm{cnt}(2d_0+4d_2,\,2d_1+2d_2)\bigr)(d)=2^{r-1+r'}.
\]
\formallabel{UpperTheory.V.Alg.M.num_holds}
\end{lemma}

\begin{lemma}[A lower bound for Euler characteristics]\label{pf:vwin-he:la}
Let $(B,R)$ be a standard setting, $M$ a graded module, $\beta\in B$ a non-zero-divisor on $M$ such that $()$ is
admitted for $\vwinheAnn(M/\beta M)$, $\delta\ge0$ in $\Z^k$, $y_1,\dots,y_n\in R_\delta$, $c_1,\dots,c_n\in\N$, and
$us$ a valid list such that $ts''=((\delta,\beta^{c_i}y_i))_{1\le i\le n}\cdot us$ is admitted for $\vwinheAnn(M)$.
Let $h(d)=\len_B(M/\beta M)_d$, let $\deg us$ be the list of degrees of $us$, and let $(\delta)^{n-1}$ be empty if
$n=0$. Then for large $d$,
\[
 \Bigl(\sum_ic_i\Bigr)\,\Delta_{((\delta)^{n-1},\,\deg us)}h(d)\ \le\ \chi(ts'',M)(d).
\]
\formallabel{UpperTheory.V.Alg.M.la_holds}
\end{lemma}

In \cref{pf:vwin-he:num}, the variables are $r$ of weight $-1$, $s$ of weight $-2$, $r'+1$ of weight $1$ and $s'+1$ of
weight $2$, so $\mathrm{cnt}(2d_0+4d_2,2d_1+2d_2)$ is a product of two counting functions of monomials of given
weighted degree, and the iterated difference computes a leading coefficient.

\begin{proposition}[Lower bound]\label{pf:vwin-he:low}
In the situation of \cref{pf:vwin-he:core}, let $a\in A_0$ with $x(a)=0$, $J=(\ker\phi)^R$, and assume that $ts$ is
admitted for $J+aR$. Then for large $d$,
\[
 2^{r+r'}\sum_{i<r}c_i\ \le\ \chi\bigl(ts,(R/J)/a\bigr)(d).
\]
\formallabel{UpperTheory.V.Alg.M.low_holds}
\end{proposition}

\begin{proof}
Let $\mathfrak q$ be as in \cref{pf:vwin-he:lk}, which applies as $F_r\in\vwinhesat(x)$; $J+aR\subset\mathfrak q$ since
$x(a)=0$. The surjection $(R/J)/a\to R/\mathfrak q$ has a kernel generated by homogeneous elements, and $ts$ is
admitted for $J+aR\subset\vwinheAnn((R/J)/a)$; by \ref{pf:vwin-he:E-gr} and \ref{pf:vwin-he:E-mono},
\begin{equation}\label{pf:vwin-he:eq-mono}
 \chi(ts,R/\mathfrak q)(d)\le\chi(ts,(R/J)/a)(d)\qquad\text{for large }d.
\end{equation}
Let $ts'=((e_0,\varpi^{2c_i}\,y_i^2t_0))_{i<r}\cdot ts_{\mathrm R}$. Since $F_i-\varpi^{c_i}y_i\in\vwinhesat(x)$, both
elements lying in $A_{-1}$, \cref{pf:vwin-he:lk} gives
$F_i^2t_0-(\varpi^{c_i}y_i)^2t_0\in\mathfrak q$, and $(\varpi^{c_i}y_i)^2t_0=\varpi^{2c_i}\,y_i^2t_0$; so $ts$ and $ts'$
agree modulo $\mathfrak q=\vwinheAnn(R/\mathfrak q)$, and $\chi(ts,R/\mathfrak q)=\chi(ts',R/\mathfrak q)$ by
\ref{pf:vwin-he:E-cong}. Moreover $ts'$ is admitted for $\mathfrak q$, since
$\vwinhesecs(ts)\subset\mathfrak q+\vwinhesecs(ts')$.

We apply \cref{pf:vwin-he:la} to $M=R/\mathfrak q$, $\beta=\varpi$, $\delta=e_0$,
$n=r$, the elements $y_i^2t_0$, the integers $2c_i$ and $us=ts_{\mathrm R}$. Here $\varpi$ is a non-zero-divisor on
$R/\mathfrak q$, and $()$ is admitted for $\vwinheAnn((R/\mathfrak q)/\varpi)$: for $b\in\mathfrak m_B$ we have
$x(b)\in\varpi\OO$ as $x$ is local, and $b-x(b)\in\mathfrak q$, so $\mathfrak m_BR\subset\mathfrak q+\varpi R$. The
degrees of $ts_{\mathrm R}$ are $(e_0)^s,(e_1)^{r'},(e_1)^{s'},e_2$, so the list of degrees in the conclusion is the
list $L$ of \cref{pf:vwin-he:num}. By \cref{pf:vwin-he:hf} and \cref{pf:vwin-he:num} (the chart has the weight
multiplicities $n_{r,s,r',s',0}$) we obtain, for large $d$,
\[
 2^{r+r'}\sum_ic_i=\Bigl(\sum_i2c_i\Bigr)2^{r-1+r'}=\Bigl(\sum_i2c_i\Bigr)\Delta_Lh(d)\le\chi(ts',R/\mathfrak q)(d)
 =\chi(ts,R/\mathfrak q)(d),
\]
and \eqref{pf:vwin-he:eq-mono} concludes the proof.
\end{proof}

\begin{proof}[Proof of \cref{pf:vwin-he:core}]
We have $a\in B=A_0$ by (M5). Let $J=(\ker\phi)^R$; it is generated by homogeneous elements (the monomials of an
element of $J$ lie in $J$), and it contains the $Z_i$, since $\phi(Z_i)=0$. By
\cref{pf:vwin-he:aux}\ref{pf:vwin-he:aux-adm}, $ts\cdot((0,a))$ is admitted for $J$; equivalently, $ts$ is admitted for
$J+aR$. The element $a$ is a non-zero-divisor on $R/J$: if $ap\in J$, then $\phi(a)\phi(p_\mu)=0$ for all $\mu$, and
$\phi(a)\ne0$ in the domain $\mathrm{Frac}(\OO)[[\tau]]$, so $p\in J$. Hence, for large $d$,
\begin{align*}
 \chi\bigl(ts,(R/J)/a\bigr)(d)&=\chi\bigl(((0,a))\cdot ts,R/J\bigr)(d)=\chi\bigl(ts\cdot((0,a)),R/J\bigr)(d)\\
 &\le\chi\bigl(((0,a)),(R/J)/\vwinhesecs(ts)\bigr)(d)=\chi\bigl(((0,a)),R/(J+\vwinhesecs(ts))\bigr)(d),
\end{align*}
by \ref{pf:vwin-he:E-nzd}, \ref{pf:vwin-he:E-perm} (as $J\subset\vwinheAnn(R/J)$), \ref{pf:vwin-he:E-ub},
\ref{pf:vwin-he:E-quot} and \ref{pf:vwin-he:E-iso}. The ideal $J'=J+\vwinhesecs(ts)$ is generated by homogeneous
elements and contains $\vwinhesecs(ts)$ and the $Z_i$. Combining \cref{pf:vwin-he:low}, this chain and
\cref{pf:vwin-he:up} for a common large $d$,
\[
 2^{r+r'}\sum_ic_i\ \le\ 2^{r+1+r'}\sum_{\xi\in T}v(\xi(a)),
\]
and dividing by $2^{r+r'}$ gives the assertion.
\end{proof}

\begin{remark}[On the formal statements]\label{pf:vwin-he:rem-M}
(1) The valuation $v$ is only required to be multiplicative and ultrametric on non-zero elements; $v(0)$ plays no
role, since the points of $T$ do not vanish at $a$.
(2) By (M2) and \cref{pf:vwin-he:vnn}, all terms $v(\xi(a))$ are non-negative; the bound is
in terms of the order of vanishing of the degree-zero function $a$ along the two-signed orbits, and (M5) requires
$a$ to vanish at $x$ but on none of these orbits.
\end{remark}

\subsection{Proofs of the auxiliary results}\label{sec:pf:vwin-he:auxproofs}

This subsection contains the proofs of the lemmas of \cref{sec:pf:vwin-he:lemmaM} whose proofs were postponed; the
notation is that of \cref{sec:pf:vwin-he:lemmaM}, and $B=A_0$. For a finite set $\tau$, a power series
$P\in\mathrm{Frac}(\OO)[[\tau]]$ and $\nu\in\N^\tau$ we write $[X^\nu]P$ for the coefficient of $X^\nu$ in $P$,
$[X_i]P$ for the coefficient of the variable $X_i$, and $\deg\nu=\sum_i\nu_i$; the \emph{linear part} of $P$ is
$\sum_i([X_i]P)X_i$, and $P$ is \emph{of pure weight $j$} if all its monomials have weight $j$. For $f\in A$ and
$j\in\Z$, $f_j\in A_j$ denotes the homogeneous component of degree $j$ of $f$.

\subsubsection{Valuations and lengths}\label{sec:pf:vwin-he:pf-val}

\begin{lemma}\label{pf:vwin-he:int}
Let $y\in K$ be non-zero and integral over $\OO$. Then $v(y)\ge0$.
\formallabel{UpperTheory.V.Alg.M.int_holds}
\end{lemma}

\begin{proof}
From $v(1)=v(1\cdot1)=2v(1)$ and $2v(-1)=v(1)$ we get $v(\pm1)=0$; hence $v(-z)=v(z)$ for $z\ne0$, and
$v(y^k)=kv(y)$. By induction on the number of terms, the ultrametric inequality shows that a non-zero finite sum of
elements of $K$, each of which is zero or has valuation at least $\kappa$, has valuation at least $\kappa$. Suppose
$v(y)<0$, and let $y^n+c_{n-1}y^{n-1}+\dots+c_0=0$ with $c_i\in\OO$; then $n\ge1$, and $S=\sum_{i<n}c_iy^i=-y^n$ is
non-zero. Every non-zero term satisfies $v(c_iy^i)=\ord_\OO(c_i)+iv(y)\ge iv(y)\ge(n-1)v(y)$, as $i\le n-1$ and
$v(y)<0$. Hence $v(S)\ge(n-1)v(y)$; but $v(S)=v(y^n)=nv(y)<(n-1)v(y)$, a contradiction.
\end{proof}

\begin{proof}[Proof of \cref{pf:vwin-he:vnn}]
The subring $\xi(A_0)$ of $K$ is the image, under the homomorphism $A/I\to K$ induced by $\xi$, of the image of
$A_0$ in $A/I$; so it is a finitely generated $\OO$-module. It contains $\OO[\xi(b)]$, which is therefore a finitely
generated $\OO$-module, $\OO$ being Noetherian. Hence $\xi(b)$ is integral over $\OO$, and \cref{pf:vwin-he:int}
applies.
\end{proof}

\begin{lemma}\label{pf:vwin-he:lqs}
Let $D$ be a domain which is an $\OO$-algebra, finitely generated as an $\OO$-module, with $\OO\to D$ injective, and
let $0\ne a\in D$. Then the set $\Hom_\OO(D,K)$ of $\OO$-algebra homomorphisms $D\to K$ is finite, $D/aD$ has finite
length over $\OO$, and
\[
 \len_\OO(D/aD)=\sum_{\sigma\in\Hom_\OO(D,K)}v(\sigma(a)).
\]
\formallabel{UpperTheory.V.Alg.M.length_quotient_span_eq_finsum_valuation}
\end{lemma}

\begin{proof}
As $D$ is a domain and $\OO\to D$ is injective, $D$ is a torsion-free, hence free, $\OO$-module of some finite rank
$n$. Multiplication by $a$ is an injective $\OO$-linear endomorphism $m_a$ of $D$ with cokernel $D/aD$. By the
elementary divisor theorem there are bases of $D$ (as source and as target) in which $m_a$ is diagonal with entries
$u_1\varpi^{c_1},\dots,u_n\varpi^{c_n}$, $u_i\in\OO^\times$; the determinant changes by a unit under these base
changes. Hence $D/aD\cong\bigoplus_i\OO/\varpi^{c_i}$ has finite length $\sum_ic_i=\ord_\OO(N(a))$, where
$N(a)=\det(m_a)\ne0$ is the norm of $a$.

Let $F=\mathrm{Frac}(\OO)\subset K$ and $L=D\otimes_\OO F$, the localization of $D$ at $\OO\setminus\{0\}$; it is a
domain of dimension $n$ over $F$, hence a field, and $N(a)=N_{L/F}(a)$ (same matrix). An $\OO$-algebra homomorphism
$D\to K$ maps $\OO\setminus\{0\}$ into $K^\times$, so it extends uniquely to an $F$-embedding $L\to K$; thus
restriction is a bijection from the set of $F$-embeddings $L\to K$ onto $\Hom_\OO(D,K)$. Since $F$ has
characteristic $0$, $L/F$ is separable, and as $K$ is algebraically closed there are exactly $n$ embeddings and
$N_{L/F}(a)=\prod_\tau\tau(a)$, the product over the $F$-embeddings $\tau:L\to K$. Hence $\Hom_\OO(D,K)$ is finite and
$N(a)=\prod_\sigma\sigma(a)$ in $K$.

Every $\sigma\in\Hom_\OO(D,K)$ is injective: $\ker\sigma$ is a prime ideal of $D$ with $\ker\sigma\cap\OO=0$, and $D$
is integral over $\OO$, so $\ker\sigma=0$ by incomparability (the prime $0\subset\ker\sigma$ lies over $0$ as well).
Hence all $\sigma(a)$ are non-zero, and as $v$ is multiplicative on non-zero elements and normalized on $\OO$,
\[
 \sum_\sigma v(\sigma(a))=v\Bigl(\prod_\sigma\sigma(a)\Bigr)=v(N(a))=\ord_\OO(N(a))=\len_\OO(D/aD).\qedhere
\]
\end{proof}

\begin{proof}[Proof of \cref{pf:vwin-he:th}]
Write $B$ for $B'$; lengths are over $B$ unless indicated otherwise. We use the following facts.
(i) $\vwinheloclen_\sigma$ is additive in short exact sequences, since localization is exact.
(ii) Let $P\subset B$ be a prime ideal. If $P=\ker\sigma$, then $(B/P)_P$ is the residue field of $B_P$ and
$\vwinheloclen_\sigma(B/P)=1$. If $P\not\subset\ker\sigma$, an element of $P\setminus\ker\sigma$ becomes a unit which
kills $(B/P)_{\ker\sigma}$, so $\vwinheloclen_\sigma(B/P)=0$.
(iii) $\len_\OO X=\len_BX$ for every $B$-module $X$, as $B$ is local with the residue field of $\OO$.
(iv) An endomorphism of a module of finite length has kernel and cokernel of the same length.
(v) Every $\OO$-algebra homomorphism $\sigma:B\to K$ is injective on $\OO$, since $\OO\to K$ is.

We show: \emph{let $I\subset B$ be an ideal such that $B/I$ is a finitely generated $\OO$-module, and $N$ a finitely
generated $B$-module with $I\subset\vwinheAnn(N)$, $\vwinheloclen_\sigma(N)<\infty$ for $\sigma\in S$ and
$\vwinheloclen_\sigma(N)=0$ for $\sigma\notin S$; then $N/aN$ and $N[a]$ have finite length and}
\begin{equation}\label{pf:vwin-he:eq-herbrand}
 \len(N/aN)-\len(N[a])=\sum_{\sigma\in S}\vwinheloclen_\sigma(N)\,v(\sigma(a)).
\end{equation}
For $I=\vwinheAnn(M)$ and $N=M$ this gives $\theta(a,M)=\sum_{\sigma\in S}\vwinheloclen_\sigma(M)\,v(\sigma(a))$,
which is at most $c\sum_{\sigma\in S}v(\sigma(a))$, as $\vwinheloclen_\sigma(M)\le c$ and $v(\sigma(a))\ge0$ for
$\sigma\in S$.

The hypotheses on $N$ pass to submodules and quotients of $N$: annihilators grow, and by (i) localized lengths do not
grow. As $B$ is Noetherian, $N$ has a finite filtration whose successive quotients are isomorphic to modules $B/P$
with $P$ prime. Hence it suffices to treat $N=0$, where both sides vanish, the modules $N\cong B/P$, and extensions.

\emph{Extensions.} Let $0\to N_1\to N_2\to N_3\to0$ be exact, with \eqref{pf:vwin-he:eq-herbrand} known for $N_1$ and
$N_3$. The snake lemma for multiplication by $a$ gives an exact sequence
\[
 0\to N_1[a]\to N_2[a]\to N_3[a]\to N_1/aN_1\to N_2/aN_2\to N_3/aN_3\to0 .
\]
Hence $N_2[a]$ and $N_2/aN_2$ have finite length, and the alternating sum of the six lengths vanishes; so the left
side of \eqref{pf:vwin-he:eq-herbrand} is additive, and by (i) so is the right side.

\emph{Prime quotients.} Let $N=B/P$ with $P$ prime. Then $I\subset P$, so $D=B/P$ is a domain and, as a quotient of
$B/I$, a finitely generated $\OO$-module; let $\pi:B\to D$ be the projection. Suppose first $P\cap\OO=0$. Then $D$ is
as in \cref{pf:vwin-he:lqs}, and every $\OO$-algebra homomorphism $\tau:D\to K$ is injective (see the proof of
\cref{pf:vwin-he:lqs}). Hence $\ker(\tau\pi)=P$ and $\vwinheloclen_{\tau\pi}(D)=1$ by (ii), so $\tau\pi\in S$. Such
$\tau$ exist, since $D\otimes_\OO\mathrm{Frac}(\OO)$ is a finite field extension of $\mathrm{Frac}(\OO)\subset K$ and
$K$ is algebraically closed; as $(\tau\pi)(a)\ne0$, we get $\pi(a)\ne0$. Since $D$ is a domain, $D[a]=0$, and by (iii)
and \cref{pf:vwin-he:lqs}
\[
 \len(D/aD)-\len(D[a])=\len_\OO(D/aD)=\sum_{\tau\in\Hom_\OO(D,K)}v(\tau(\pi(a))).
\]
On the right side of \eqref{pf:vwin-he:eq-herbrand}, an element $\sigma\in S$ contributes $0$ unless
$P\subset\ker\sigma$, by (ii); in that case $\sigma=\tau\pi$ for a unique $\tau$, and $\vwinheloclen_\sigma(D)=1$.
Since $\tau\mapsto\tau\pi$ is injective with image in $S$, both sides agree. Suppose now $P\cap\OO\ne0$. Then
$P\cap\OO$ is the maximal ideal of $\OO$, and since $D$ is integral over $\OO$, $P$ is maximal; so $D$ is a simple
$B$-module, and the left side of \eqref{pf:vwin-he:eq-herbrand} vanishes by (iv). If $P\subset\ker\sigma$, then
$\sigma$ kills a non-zero element of $\OO$, contradicting (v); so $\vwinheloclen_\sigma(D)=0$ for all $\sigma$ by
(ii), and the right side vanishes as well.
\end{proof}

\subsubsection{Charts}\label{sec:pf:vwin-he:pf-chart}

\begin{lemma}\label{pf:vwin-he:lin}
Let $(A,x)$ be a weighted ring and $\phi:A\to\mathrm{Frac}(\OO)[[\tau]]$ an $\OO$-algebra homomorphism such that the
constant term of $\phi(f)$ is $x(f)$ for all $f\in A$.
\begin{enumerate}[label=(\alph*)]
\item\label{pf:vwin-he:lin-sat} For $f\in\mathfrak m_x^2$, $\phi(f)$ has no terms of degree $\le1$. In particular
$[X_i]\phi(f)=0$ for all $f\in\vwinhesat(x)$ and $i\in\tau$.
\item\label{pf:vwin-he:lin-wt} If moreover $\phi(A_j)$ consists of series of pure weight $j$ for every $j$, then
$[X_i]\phi(f)=[X_i]\phi(f_{\mathrm{wt}(i)})$ for all $f\in A$, and $[X_i]\phi(f)=0$ for $f\in A_j$ with
$j\ne\mathrm{wt}(i)$.
\end{enumerate}
\end{lemma}

\begin{proof}
(a) For $f,g\in\mathfrak m_x$ the series $\phi(f)$, $\phi(g)$ have no constant term, so $\phi(fg)=\phi(f)\phi(g)$
has no terms of degree $\le1$, and this passes to sums. If $cf\in\mathfrak m_x^2$ with $0\ne c\in\OO$, then
$c\,[X_i]\phi(f)=[X_i]\phi(cf)=0$, and $c$ is invertible in $\mathrm{Frac}(\OO)$. (b) The variable $X_i$ has weight
$\mathrm{wt}(i)$, so it does not occur in a series of pure weight $j\ne\mathrm{wt}(i)$; apply this to the components
of $f$.
\end{proof}

\begin{proof}[Proof of \cref{pf:vwin-he:rk}]
Let $\phi:A\to\mathrm{Frac}(\OO)[[\tau]]$ be a chart with $\#\{i:\mathrm{wt}(i)=j\}=n(j)$ for all $j$, and let
$\tau_{-1}$ be the set of variables of weight $-1$. As $x(A_{-1})=0$, $\vwinhecot_{-1}(x)$ is the image of $A_{-1}$ in
$A/\vwinhesat(x)$. By \cref{pf:vwin-he:fg}, $A_{-1}$ is generated over $A_0$ by finitely many elements
$s_1,\dots,s_k$. For $b\in A_0$ we have $bs_j-x(b)s_j=(b-x(b))s_j\in\mathfrak m_x^2$, since $b-x(b)$ and $s_j$ lie in
$\mathfrak m_x$; so the images of $\sum_jb_js_j$ and of $\sum_jx(b_j)s_j$ in $A/\vwinhesat(x)$ coincide, and the
images of the $s_j$ generate $\vwinhecot_{-1}(x)$ over $\OO$.

By \cref{pf:vwin-he:lin}\ref{pf:vwin-he:lin-sat}, $f\mapsto([X_i]\phi(f))_{i\in\tau_{-1}}$ induces an $\OO$-linear map
$\Lambda:A/\vwinhesat(x)\to\mathrm{Frac}(\OO)^{\tau_{-1}}$. It is injective on $\vwinhecot_{-1}(x)$: if $f\in A_{-1}$
and $\Lambda(f)=0$, then also $[X_i]\phi(f)=0$ for $\mathrm{wt}(i)\ne-1$ by
\cref{pf:vwin-he:lin}\ref{pf:vwin-he:lin-wt}, so $\phi(f)$ has no linear terms, and $f\in\vwinhesat(x)$ by the last
property of charts. Hence $\vwinhecot_{-1}(x)$ is torsion-free, thus free, of rank at most $\#\tau_{-1}$, since
$\OO$-linearly independent elements of $\mathrm{Frac}(\OO)^{\tau_{-1}}$ are linearly independent over
$\mathrm{Frac}(\OO)$. Conversely, for $i\in\tau_{-1}$ choose $f^{(i)}\in A$ whose linear part is $u_iX_i$ with
$u_i\ne0$, and let $g^{(i)}=f^{(i)}_{-1}$. By \cref{pf:vwin-he:lin}\ref{pf:vwin-he:lin-wt}, $\Lambda(g^{(i)})$ is $u_i$
times the $i$-th unit vector, so the images of the $g^{(i)}$ in $\vwinhecot_{-1}(x)$ are linearly independent, and the
rank is at least $\#\tau_{-1}=n(-1)$.
\end{proof}

\begin{lemma}[Formal implicit functions]\label{pf:vwin-he:ift}
Let $E$ be a field, $\tau$ a finite set with a map $\mathrm{wt}$ to an abelian group $Y$, $\tau_0$ the set of
$i\in\tau$ with $\mathrm{wt}(i)=0$ and $\tau'=\tau\setminus\tau_0$ with the restricted weights. Let $z_k\in E[[\tau]]$,
$k\in\iota$, be finitely many series of pure weight $0$ without constant term, where $\#\iota=\#\tau_0$, and assume
that for every $i\in\tau_0$ there are $c_{ik}\in E$ such that $\sum_kc_{ik}z_k$ has linear part $X_i$. Then there is an
$E$-algebra homomorphism $\psi:E[[\tau]]\to E[[\tau']]$ which maps series of pure weight $j$ to series of pure weight
$j$ for every $j\in Y$, such that $\psi(z_k)=0$ for all $k$, the constant term of $\psi(g)$ is that of $g$, and
$[X_i]\psi(g)=[X_i]g$ for all $g\in E[[\tau]]$ and $i\in\tau'$.
\formallabel{UpperTheory.V.Alg.M.ift_holds}
\end{lemma}

\begin{proof}
For a family $h=(h_i)_{i\in\tau}$ of series in $E[[\tau']]$ without constant term, let $h^*:E[[\tau]]\to E[[\tau']]$
be the substitution $X_i\mapsto h_i$; it is a well-defined $E$-algebra homomorphism, as a monomial $X^\nu$ is mapped
to a series without terms of degree $<\deg\nu$. If each $h_i$ is of pure weight $\mathrm{wt}(i)$, then $h^*$ preserves
pure weights. Put $G_i=X_i$ for $i\in\tau'$ and $G_i=0$ for $i\in\tau_0$, and $g_i=\sum_kc_{ik}z_k-X_i$ for
$i\in\tau_0$ and $g_i=0$ for $i\in\tau'$. Each $g_i$ has no terms of degree $\le1$ and is of pure weight
$\mathrm{wt}(i)$. We construct $h$ with
\begin{equation}\label{pf:vwin-he:eq-ift}
 h_i=G_i-h^*(g_i)\qquad(i\in\tau).
\end{equation}
Let $h^{(0)}=0$ and $h^{(n+1)}_i=G_i-(h^{(n)})^*(g_i)$. By induction, each $h^{(n)}_i$ has no constant term and is of
pure weight $\mathrm{wt}(i)$. If two families $h,h'$ agree in degrees $<n$, then $h^*(g)$ and $h'^*(g)$ agree in
degrees $<n+1$ for every $g$ without terms of degree $\le1$, since each monomial of such a $g$ is a product of at
least two variables and is mapped to a product of at least two series without constant term. Hence $h^{(n+1)}$ and
$h^{(n)}$ agree in degrees $<n+1$ for all $n$, the $h^{(n)}$ converge coefficientwise to a family $h$ without constant
terms, with $h_i$ of pure weight $\mathrm{wt}(i)$, and $h$ satisfies \eqref{pf:vwin-he:eq-ift}. Let $\psi=h^*$.

Then $\psi$ preserves pure weights, and constant terms since the $h_i$ have none. By \eqref{pf:vwin-he:eq-ift},
$h_i-G_i$ has no terms of degree $\le1$; hence $\psi(g)$ is congruent to $g(0)+\sum_i([X_i]g)G_i$ modulo terms of
degree $\ge2$, and $[X_i]\psi(g)=[X_i]g$ for $i\in\tau'$. Finally, for $i\in\tau_0$, \eqref{pf:vwin-he:eq-ift} reads
$h_i=-\sum_kc_{ik}\psi(z_k)+\psi(X_i)=-\sum_kc_{ik}\psi(z_k)+h_i$, so $\sum_kc_{ik}\psi(z_k)=0$. The square matrices
$C=(c_{ik})_{i\in\tau_0,k\in\iota}$ and $\Lambda=([X_i]z_k)_{k\in\iota,i\in\tau_0}$ satisfy $C\Lambda=1$, because
$\sum_kc_{ik}[X_{i'}]z_k=[X_{i'}]X_i$ for $i,i'\in\tau_0$; hence $\Lambda C=1$, and
$\psi(z_k)=\sum_{i}\Lambda_{ki}\sum_{k'}c_{ik'}\psi(z_{k'})=0$.
\end{proof}

\begin{proof}[Proof of \cref{pf:vwin-he:cut}]
By (M1) there is a chart $\phi_1:A\to\mathrm{Frac}(\OO)[[\tau_1]]$ with $\#\{i:\mathrm{wt}(i)=j\}=n_{r,s,r',s',m}(j)$
for all $j$; the set $\tau_0\subset\tau_1$ of variables of weight $0$ has $m$ elements. Let $i\in\tau_0$. Choose
$f\in A$ such that the linear part of $\phi_1(f)$ is $uX_i$ with $u\ne0$, and put $g=f_0-x(f_0)\in A_0\cap\mathfrak
m_x$. By \cref{pf:vwin-he:lin}\ref{pf:vwin-he:lin-wt} the linear part of $\phi_1(f_0)$, and hence that of
$\phi_1(g)$, is again $uX_i$. Since the $Z_k$ span the cotangent space in degree $0$, there are $0\ne c\in\OO$ and
$a_k\in\OO$ with $\sum_ka_kZ_k-cg\in\vwinhesat(x)$; by \cref{pf:vwin-he:lin}\ref{pf:vwin-he:lin-sat}, the series
$\sum_k(a_k/cu)\phi_1(Z_k)$ has linear part $X_i$. The series $z_k=\phi_1(Z_k)$ are of pure weight $0$ without
constant term, as $Z_k\in A_0\cap\mathfrak m_x$. So \cref{pf:vwin-he:ift} (with $E=\mathrm{Frac}(\OO)$ and $Y=\Z$)
gives $\psi:\mathrm{Frac}(\OO)[[\tau_1]]\to\mathrm{Frac}(\OO)[[\tau]]$, where $\tau=\tau_1\setminus\tau_0$. Let
$\phi=\psi\circ\phi_1$. Then $\phi(Z_k)=0$; $\phi(A_j)$ consists of series of pure weight $j$; the constant term of
$\phi(f)$ is $x(f)$; and for $i\in\tau$ and $f$ such that the linear part of $\phi_1(f)$ is $uX_i$, the linear part of
$\phi(f)$ is the restriction of that of $\phi_1(f)$ to the variables in $\tau$, i.e.\ $uX_i$. Every $i\in\tau$ has
$\mathrm{wt}(i)\ne0$ and $n_{r,s,r',s',m}(\mathrm{wt}(i))\ge1$, so $\mathrm{wt}(i)\in\{-2,-1,1,2\}$. Thus $\phi$ is a
reduced chart, and $\#\{i\in\tau:\mathrm{wt}(i)=j\}$ equals $n_{r,s,r',s',m}(j)=n_{r,s,r',s',0}(j)$ for $j\ne0$, and
$0=n_{r,s,r',s',0}(0)$ for $j=0$.
\end{proof}

\begin{proof}[Proof of \cref{pf:vwin-he:nz}]
If $\phi(a)\ne0$, take $a'=a$. Otherwise choose variables $i$ of weight $-1$ and $i'$ of weight $1$. By the third
property of charts and \cref{pf:vwin-he:lin}\ref{pf:vwin-he:lin-wt} (applied to a homogeneous component), there are
$f\in A_{-1}$ with $[X_i]\phi(f)\ne0$ and $g\in A_1$ with $[X_{i'}]\phi(g)\ne0$; in particular $\phi(g)\ne0$. We
claim that $\phi(F_k)\ne0$ for some $k$. Otherwise the $\OO$-linear map $\ell:A\to\mathrm{Frac}(\OO)$,
$\ell(y)=[X_i]\phi(y)$, vanishes on all $F_k$ and, by \cref{pf:vwin-he:lin}\ref{pf:vwin-he:lin-sat}, on
$\vwinhesat(x)$. As $f\in A_{-1}\subset\mathfrak m_x$ and the $F_k$ span the cotangent space in degree $-1$, there is
$0\ne c\in\OO$ with $cf\in\sum_k\OO F_k+\vwinhesat(x)$; so $c\,\ell(f)=0$ and $\ell(f)=0$, a contradiction. Now put
$a'=a+gF_k$. Then $gF_k\in A_1A_{-1}\subset A_0$ and $gF_k\in I$; $x(a')=x(a)+x(g)x(F_k)=x(a)$, as $x(A_1)=0$; and
$\phi(a')=\phi(g)\phi(F_k)\ne0$, since $\mathrm{Frac}(\OO)[[\tau]]$ is a domain.
\end{proof}

\begin{lemma}\label{pf:vwin-he:filt}
Let $(A,x)$ be a weighted ring and $\phi:A\to\mathrm{Frac}(\OO)[[\tau]]$ a reduced chart.
\begin{enumerate}[label=(\alph*)]
\item\label{pf:vwin-he:filt-mul} If $P\in\mathrm{Fil}^a$ and $Q\in\mathrm{Fil}^b$, then $PQ\in\mathrm{Fil}^{a+b}$, and
a series of pure weight $j$ lies in $\mathrm{Fil}^{\max(-j,0)}$. (This holds for power series over any commutative
ring and for any weights.)
\item\label{pf:vwin-he:filt-sat} If $f\in A_{-1}\cap\vwinhesat(x)$, then $\phi(f)\in\mathrm{Fil}^2$; if $b\in A_0$ and
$x(b)=0$, then $\phi(b)\in\mathrm{Fil}^1$.
\item\label{pf:vwin-he:filt-lat} Let $F\in A_{-1}\cap\vwinhesat(x)$, and let $N^\mu$ be formed as in
\cref{pf:vwin-he:rees} with $F$ in place of $F_r$. Then $N^\mu\subset A_{2(\mu_1-\mu_0-\mu_2)}$ and
$\phi(N^\mu)\subset\mathrm{Fil}^{2\mu_0+4\mu_2}$ for all $\mu\in\N^3$.
\end{enumerate}
\formallabel{UpperTheory.V.Alg.M.fil_holds}
\formallabel{UpperTheory.V.Alg.M.satfil_holds}
\formallabel{UpperTheory.V.Alg.M.latfil_holds}
\end{lemma}

\begin{proof}
(a) The coefficient of $X^\nu$ in $PQ$ is $\sum_{\nu'+\nu''=\nu}([X^{\nu'}]P)([X^{\nu''}]Q)$, and
$|\nu|_-=|\nu'|_-+|\nu''|_-$; if $|\nu|_-<a+b$, then in each term $|\nu'|_-<a$ or $|\nu''|_-<b$. A monomial $\nu$ of
weight $j=|\nu|_+-|\nu|_-$ has $|\nu|_-=|\nu|_+-j\ge-j$.

(b) As the variables have weights in $\{-2,-1,1,2\}$, every monomial satisfies $\deg\nu\le|\nu|_-+|\nu|_+$. Let
$f\in A_{-1}\cap\vwinhesat(x)$, and let $\nu$ satisfy $|\nu|_-<2$ and $[X^\nu]\phi(f)\ne0$. Then $\nu$ has weight
$-1$, so $|\nu|_+=|\nu|_--1$, whence $|\nu|_-=1$, $|\nu|_+=0$ and $\deg\nu\le1$. But $cf\in\mathfrak m_x^2$ for some
$0\ne c\in\OO$, so $\phi(cf)=c\phi(f)$ has no terms of degree $\le1$ by
\cref{pf:vwin-he:lin}\ref{pf:vwin-he:lin-sat}, a contradiction. Similarly, if $b\in A_0$, $x(b)=0$, $|\nu|_-=0$ and
$[X^\nu]\phi(b)\ne0$, then $\nu$ has weight $0$, so $|\nu|_+=0$ and $\nu=0$; but the constant term of $\phi(b)$ is
$x(b)=0$.

(c) For $j\in\Z$ and $n\in\N$ let $S(j,n)=\{f\in A_j:\phi(f)\in\mathrm{Fil}^n\}$. By (a), and as
$\mathrm{Fil}^0$ contains all series, these are $A_0$-submodules with $S(j,n)S(j',n')\subset S(j+j',n+n')$, and
$A_{-2}=S(-2,2)$, $A_2=S(2,0)$. By (b) and (a), $\phi(F^2)\in\mathrm{Fil}^4$, so $BF^2\subset S(-2,4)$; also
$A_{-2}A_{-2}A_2\subset S(-2,4)$. Hence $N_0\subset S(-2,2)$, $N_1\subset S(2,0)$, $N_2\subset S(-2,4)$, and
$N^\mu\subset S(2\mu_1-2\mu_0-2\mu_2,\,2\mu_0+4\mu_2)$.
\end{proof}

\begin{proof}[Proof of \cref{pf:vwin-he:lk}]
Write $n(\mu)=2\mu_0+4\mu_2$. Clearly $\mathfrak q$ is an additive subgroup of $R$. The ring $R$ is spanned by the
elements $ut^{\mu'}$ with $u\in N^{\mu'}$, and for $p\in\mathfrak q$ the coefficient of $t^\mu$ in $ut^{\mu'}p$ is
$up_{\mu-\mu'}$ if $\mu\ge\mu'$ and $0$ otherwise. By \cref{pf:vwin-he:filt}\ref{pf:vwin-he:filt-lat} (with $F=F_r$)
and \ref{pf:vwin-he:filt-mul}, $\phi(u)\in\mathrm{Fil}^{n(\mu')}$ and
$\phi(up_{\mu-\mu'})\in\mathrm{Fil}^{n(\mu')+n(\mu-\mu')+1}=\mathrm{Fil}^{n(\mu)+1}$; so $\mathfrak q$ is an ideal.
Membership in $\mathfrak q$ is tested coefficientwise, so the homogeneous components $p_\mu t^\mu$ of an element
$p\in\mathfrak q$ lie in $\mathfrak q$, and $\mathfrak q$ is generated by homogeneous elements. If all coefficients of
$p$ lie in $\ker\phi$, then clearly $p\in\mathfrak q$. For $b\in B$ with $x(b)=0$, the only coefficient of
$b\in R_0$ is $b$, and $\phi(b)\in\mathrm{Fil}^1$ by \cref{pf:vwin-he:filt}\ref{pf:vwin-he:filt-sat}. If
$cp\in\mathfrak q$ with $0\ne c\in\OO$, then $c\phi(p_\mu)\in\mathrm{Fil}^{n(\mu)+1}$ for all $\mu$, and so is
$\phi(p_\mu)$, as $\mathrm{Fil}^n$ is a $\mathrm{Frac}(\OO)$-subspace. Finally, for $f,g\in A_{-1}$ with
$f-g\in\vwinhesat(x)$, the element $f^2t_0-g^2t_0$ has the single coefficient $f^2-g^2=(f-g)(f+g)$, at $\mu=e_0$; here
$\phi(f-g)\in\mathrm{Fil}^2$ by \cref{pf:vwin-he:filt}\ref{pf:vwin-he:filt-sat} and $\phi(f+g)\in\mathrm{Fil}^1$ by
\cref{pf:vwin-he:filt}\ref{pf:vwin-he:filt-mul}, as $f+g$ has pure weight $-1$; so
$\phi(f^2-g^2)\in\mathrm{Fil}^3=\mathrm{Fil}^{n(e_0)+1}$.
\end{proof}

\begin{lemma}\label{pf:vwin-he:diag}
Let $B$ be an $\OO$-algebra, $x:B\to\OO$ an $\OO$-algebra homomorphism, $L$ a finitely generated $B$-module, $I$ a
finite set and $\Lambda:L\to\mathrm{Frac}(\OO)^I$ an $\OO$-linear map with $\Lambda(bf)=x(b)\Lambda(f)$ for $b\in B$
and $f\in L$, such that for every $i\in I$ some $f\in L$ has $\Lambda(f)_i\ne0$ and $\Lambda(f)_{i'}=0$ for all
$i'\ne i$. Let $\pi:L\to P$ be a surjective $B$-linear map such that $\pi(f)=0$ if and only if
$\Lambda(f)=\varpi\Lambda(g)$ for some $g\in L$. Then $\len_BP=\#I$.
\end{lemma}

\begin{proof}
If $f_1,\dots,f_k$ generate $L$ over $B$, then $\Lambda(\sum_jb_jf_j)=\sum_jx(b_j)\Lambda(f_j)$; so $\Lambda(L)$ is
the $\OO$-span of the $\Lambda(f_j)$, a finitely generated torsion-free, hence free, $\OO$-submodule of
$\mathrm{Frac}(\OO)^I$. Its rank is at most $\#I$, and at least $\#I$ by the hypothesis on the coordinates; so
$\Lambda(L)\cong\OO^I$ and $\Lambda(L)/\varpi\Lambda(L)\cong(\OO/\varpi)^I$. By hypothesis the kernel of the
surjection $L\to\Lambda(L)/\varpi\Lambda(L)$ is $\ker\pi$, so $P\cong(\OO/\varpi)^I$ as $\OO$-modules. Moreover $B$
acts on $P$ through $x$: $\pi(bf)-x(b)\pi(f)=\pi(bf-x(b)f)=0$, as $\Lambda(bf-x(b)f)=0$. Hence the $B$-submodules of
$P$ are its $\OO$-submodules, and $\len_BP=\len_\OO P=\#I$.
\end{proof}

\begin{proof}[Proof of \cref{pf:vwin-he:hf}]
Let $\mu=d\in\N^3$, $n=2\mu_0+4\mu_2$, $w=2(\mu_1-\mu_0-\mu_2)$, and let $P$ be the image of $R_\mu=N^\mu t^\mu$ in
$R/(\mathfrak q+\varpi R)$, which is the piece of degree $\mu$ of $(R/\mathfrak q)/\varpi(R/\mathfrak q)$. Let
$\pi:N^\mu\to P$ map $f$ to the class of $ft^\mu$; it is $B$-linear and surjective. As $\mathfrak q$ is defined
coefficientwise, $ft^\mu\in\mathfrak q+\varpi R$ if and only if $\phi(f-\varpi g)\in\mathrm{Fil}^{n+1}$ for some
$g\in N^\mu$ (if $ft^\mu-\varpi p\in\mathfrak q$, take $g=p_\mu$).

Let $I$ be the set of $\nu\in\N^\tau$ with $|\nu|_-=n$ and $|\nu|_+=2\mu_1+2\mu_2$; it is finite, with
$\mathrm{cnt}(2d_0+4d_2,2d_1+2d_2)$ elements. Let $\Lambda:N^\mu\to\mathrm{Frac}(\OO)^I$ be
$\Lambda(f)=([X^\nu]\phi(f))_{\nu\in I}$. By \cref{pf:vwin-he:filt}\ref{pf:vwin-he:filt-lat} (with $F=F_r$), every
$f\in N^\mu$ lies in $A_w$ and satisfies $\phi(f)\in\mathrm{Fil}^n$. Hence the monomials of negative weight $<n+1$
which can occur in $\phi(f)$ have negative weight $n$ and weight $w$, so positive weight $w+n=2\mu_1+2\mu_2$; that is,
they lie in $I$, and
\begin{equation}\label{pf:vwin-he:eq-lead}
 \phi(f)\in\mathrm{Fil}^{n+1}\iff\Lambda(f)=0\qquad(f\in N^\mu).
\end{equation}

We check the hypotheses of \cref{pf:vwin-he:diag} for $L=N^\mu$, which is finitely generated over $B$ by
\cref{pf:vwin-he:fg}, and for $\Lambda$ and $\pi$. For $b\in B$, the element $b-x(b)\in A_0$ is killed by $x$, so
$\phi(b-x(b))\in\mathrm{Fil}^1$ by \cref{pf:vwin-he:filt}\ref{pf:vwin-he:filt-sat}, and
$\phi((b-x(b))f)\in\mathrm{Fil}^{n+1}$ for $f\in N^\mu$; by \eqref{pf:vwin-he:eq-lead}, $\Lambda(bf)=x(b)\Lambda(f)$.
Next we use a separation property of reduced charts (\lref{UpperTheory.V.Alg.M.exists_mem_oz_pow_coeff_ne_zero_unique}):
for $k,l\in\N$, every monomial $\nu$ with $|\nu|_-=2k$ and $|\nu|_+=2l$ admits $f\in(A_{-2})^k(A_2)^l$ with
$[X^\nu]\phi(f)\ne0$ and $[X^{\nu'}]\phi(f)=0$ for all other $\nu'$ with $|\nu'|_-=2k$. For $\nu\in I$ we apply it with
$k=\mu_0+2\mu_2$ and $l=\mu_1+\mu_2$; the resulting $f$ lies in
$(A_{-2})^{\mu_0}(A_2)^{\mu_1}(A_{-2}A_{-2}A_2)^{\mu_2}\subset N^\mu$, and $\Lambda(f)$ is supported exactly at $\nu$.
Finally, by the first paragraph and \eqref{pf:vwin-he:eq-lead}, $\pi(f)=0$ if and only if
$\Lambda(f)=\varpi\Lambda(g)$ for some $g\in N^\mu$. Hence $\len_BP=\#I$.
\end{proof}

\subsubsection{Points, slices and localized lengths}\label{sec:pf:vwin-he:pf-points}

We use the following finiteness theorem of the formal development, which we call the \emph{slice theorem}
(\lref{UpperTheory.V.Alg.M.slice_holds}): if $(A,x)$ is a weighted ring, $J\subset A$ a homogeneous ideal such
that the image of $A_0$ in $A/J$ is a finitely generated $\OO$-module and \textup{(M4)} holds for $J$, and
$f\in A_{-2}$, $g\in A_2$, then the set of $K$-points $\xi$ on $J$ with $\xi(f)=1$ and $\xi(g)\ne0$ is finite, and for
every $h\in A$ vanishing at all these points there are $n\in\N$ and $0\ne c\in\OO$ with $cg^nh\in J+(f-1)A$.

\begin{proof}[Proof of \cref{pf:vwin-he:fin}]
By \cref{pf:vwin-he:fg}, $A_{-2}$ and $A_2$ are generated over $A_0$ by finite sets $s_-$ and $s_+$. Let $T$ be the
set of two-signed $K$-points $\xi$ on $J$ such that $\xi(f)=1$ and $\xi(g)\ne0$ for some $f\in s_-$ and $g\in s_+$; it
is finite by the slice theorem. Let $\xi$ be a two-signed $K$-point on $J$. By \cref{pf:vwin-he:sign} there are
$f'\in A_{-2}$ and $g'\in A_2$ with $\xi(f')\ne0\ne\xi(g')$; writing $f'$ and $g'$ as $A_0$-linear combinations of
$s_-$ and $s_+$, we find $f\in s_-$ and $g\in s_+$ with $\xi(f)\ne0\ne\xi(g)$. Choose $t\in K^\times$ with
$t^2=\xi(f)$ and let $\xi'=t\cdot\xi$. Then $\xi'(f)=t^{-2}\xi(f)=1$ and $\xi'(g)=t^2\xi(g)\ne0$; $\xi'$ vanishes on
$J$, as $J$ is homogeneous and $\xi'(h)=\sum_jt^j\xi(h_j)$; and $\xi'$ is two-signed by \cref{pf:vwin-he:sign}. So
$\xi'\in T$, and $\xi=t^{-1}\cdot\xi'$.
\end{proof}

\begin{lemma}\label{pf:vwin-he:deh}
Let $A=\bigoplus_jA_j$ be a $\Z$-graded $\OO$-algebra, $J\subset A$ a homogeneous ideal, $f\in A_d$ with $d\ne0$, and
$h\in A_w$ with $h\in J+(f-1)A$. Then $f^nh\in J$ for some $n\in\N$.
\formallabel{UpperTheory.V.Alg.M.deh_holds}
\end{lemma}

\begin{proof}
Write $h=y+u(f-1)$ with $y\in J$ and $u\in A$, and denote classes in $A/J$ by a bar. Comparing the components of
degree $j+d$, and using $y_{j+d}\in J$, gives $\bar h_{j+d}=\bar u_j\bar f-\bar u_{j+d}$ for all $j$, where
$h_{j+d}=0$ unless $j+d=w$. Hence $\bar u_{j+d}=\bar u_j\bar f$ for $j+d\ne w$, and
$\bar u_w=\bar u_{w-d}\bar f-\bar h$. Iterating along the progression $w+\Z d$ gives $\bar u_{w+kd}=\bar u_w\bar f^k$
and $\bar u_{w-d}=\bar u_{w-kd}\bar f^{k-1}$ for $k\ge1$. As only finitely many $u_j$ are non-zero and $d\ne0$, for
large $k$ we get $\bar u_w\bar f^k=0$ and $\bar u_{w-d}=0$; hence $\bar u_w=-\bar h$ and $\bar f^k\bar h=0$, i.e.\
$f^kh\in J$.
\end{proof}

\begin{lemma}\label{pf:vwin-he:nil}
Let $A=\bigoplus_jA_j$ be a $\Z$-graded $\OO$-algebra, $J\subset A$ an ideal such that the image of $B=A_0$ in $A/J$
is a finitely generated $\OO$-module, $\sigma:B\to K$ an $\OO$-algebra homomorphism, and $p\in B$ with $\sigma(p)=0$.
Then there are $b\in B$ with $\sigma(b)\ne0$ and $n\in\N$ with $bp^n\in J$.
\end{lemma}

\begin{proof}
The image of $p$ in $A/J$ lies in the image of $B$, which is a ring and a finitely generated $\OO$-module; so it is
integral over $\OO$, and $Q(p)\in J$ for some monic $Q\in\OO[X]$. Write $Q=X^nQ_1$ with $Q_1(0)\ne0$ and put
$b=Q_1(p)\in B$. Then $bp^n=Q(p)\in J$, and $\sigma(b)=Q_1(\sigma(p))=Q_1(0)\ne0$ in $K$, as $\OO\to K$ is injective.
\end{proof}

\begin{lemma}\label{pf:vwin-he:loclen}
Let $B$ be an $\OO$-algebra, $\sigma:B\to K$ an $\OO$-algebra homomorphism and $P=\ker\sigma$, a prime ideal.
\begin{enumerate}[label=(\alph*)]
\item\label{pf:vwin-he:loclen-surj} Let $\varphi:L\to M$ be a surjective and $\psi:L\to N$ an arbitrary homomorphism
of $B$-modules such that for every $y\in\ker\psi$ there is $t\in B$ with $\sigma(t)\ne0$ and $t\varphi(y)=0$. Then
$\vwinheloclen_\sigma(M)\le\vwinheloclen_\sigma(N)$.
\item\label{pf:vwin-he:loclen-one} Let $M$ be a $B$-module and $m_0\in M$ such that for all $p\in P$ and $m\in M$
there is $b\notin P$ with $bpm=0$, and for every $m\in M$ there are $b\notin P$ and $a\in B$ with $bm=am_0$. Then
$\vwinheloclen_\sigma(M)\le1$.
\item\label{pf:vwin-he:loclen-zero} Let $M$ be a finitely generated $B$-module, $\mathfrak a,\mathfrak b\subset B$
ideals and $n\in\N$ with $\mathfrak a^nM\subset\mathfrak bM$ and $\mathfrak b\subset P$, and assume $\sigma(b)\ne0$ for
some $b\in\mathfrak a$. Then $\vwinheloclen_\sigma(M)=0$.
\end{enumerate}
\end{lemma}

\begin{proof}
(a) Localization at $P$ is exact, so $\varphi_P$ is surjective; and $\ker\psi_P\subset\ker\varphi_P$, for if
$\psi(y)/s=0$, then $\psi(uy)=0$ for some $u\notin P$, and $t\varphi(uy)=0$ with $tu\notin P$. Hence $M_P$ is a
quotient of $L_P/\ker\psi_P$, which is isomorphic to $\psi_P(L_P)\subset N_P$. (b) The map $B_P\to M_P$,
$\rho\mapsto\rho m_0$, is surjective, as $m/s=(a/bs)m_0$, and it kills $PB_P$, since $pm_0=0$ in $M_P$ for $p\in P$;
so $M_P$ is a quotient of the residue field of $B_P$. (c) If $M_P\ne0$, then $P$ lies in the support of
$M/\mathfrak bM$, as $M$ is finitely generated and $\mathfrak b\subset P$; but $b^n$ kills $M/\mathfrak bM$, so
$b^n\in P$, contradicting $\sigma(b)\ne0$.
\end{proof}

\begin{proof}[Proof of \cref{pf:vwin-he:loc3}]
Let $P=\ker\sigma$. As $\sigma$ does not vanish on $I_D$, there are $f\in A_{-2}$ and $g\in A_2$ such that
$b_0=fg\in B$ satisfies $\sigma(b_0)\ne0$.

\emph{Step 1: if $h\in A_j$, $b\in B$, $\sigma(b)\ne0$ and $bh^n\in J$ for some $n$, then $b'h\in J$ for some
$b'\in B$ with $\sigma(b')\ne0$.} Every $K$-point $\xi$ on $J$ with $\xi(f)=1$ and $\xi(g)\ne0$ satisfies
$\xi(b)\xi(h)^n=0$, hence $\xi(bh)=0$. By the slice theorem, $cg^Nbh\in J+(f-1)A$ for some $N\in\N$ and
$0\ne c\in\OO$. The element $cg^Nbh$ is homogeneous, so \cref{pf:vwin-he:deh} gives $m$ with $f^mcg^Nbh\in J$. Then
$b'=cb_0^{N+m}b\in B$ satisfies $\sigma(b')=c\,\sigma(b_0)^{N+m}\sigma(b)\ne0$ and
$b'h=f^Ng^m\cdot f^mcg^Nbh\in J$.

\emph{Step 2: for $p\in P$ and $h\in A_w$ there is $b\in B$ with $\sigma(b)\ne0$ and $bph\in J$.} By
\cref{pf:vwin-he:nil} there are $b$ with $\sigma(b)\ne0$ and $n$ with $bp^n\in J$; then $b(ph)^n\in J$, and Step 1
applies to $ph\in A_w$.

\emph{Step 3: there is $h_0\in A_w$ such that for every $h\in A_w$ there are $b,a\in B$ with $\sigma(b)\ne0$ and
$bh-ah_0\in J$.} Let $e=f^k$ and $e'=g^k$ if $w=k\ge0$, and $e=g^k$ and $e'=f^k$ if $w=-k<0$. Then $e\in A_{-2w}$ and
$ee'=b_0^k$, so $hh'e\in B$ for $h,h'\in A_w$. If $\sigma(h_0h_0e)\ne0$ for some $h_0\in A_w$, take $b=h_0h_0e$ and
$a=hh_0e$, so that $bh-ah_0=0$. Otherwise let $h_0=0$ and $h\in A_w$. As $\sigma(hhe)=0$, \cref{pf:vwin-he:nil} gives
$b$ with $\sigma(b)\ne0$ and $n$ with $b(hhe)^n\in J$; then $bb_0^{kn}h^{2n}=e'^n\,b(hhe)^n\in J$ and
$\sigma(bb_0^{kn})\ne0$, and Step 1 gives $b'h\in J$ with $\sigma(b')\ne0$.

By Steps 2 and 3, \cref{pf:vwin-he:loclen}\ref{pf:vwin-he:loclen-one} applies to $M=(A/J)_w$ and the class $m_0$ of
$h_0$, and $\vwinheloclen_\sigma((A/J)_w)\le1$.
\end{proof}

\subsubsection{The Rees algebra away from the fixed point}\label{sec:pf:vwin-he:pf-rees}

\begin{proof}[Proof of \cref{pf:vwin-he:loc1}]
Choose $f\in A_{-2}$ and $g\in A_2$ such that $b_0=fg$ satisfies $\sigma(b_0)\ne0$. Write $d=\mu\in\N^3$ and
$w(\nu)=2(\nu_1-\nu_0-\nu_2)$ for $\nu\in\N^3$. As $N_0,N_2\subset A_{-2}$ (recall $F_r^2\in A_{-2}$) and $N_1=A_2$, we
have $N^\nu\subset A_{w(\nu)}$.

(i) \emph{For $\nu\in\N^3$ and $c\in A_{w(\nu)}$ there is $n$ with $b_0^nc\in N^\nu$.} Let $k=\nu_0+2\nu_2$ and
$l=\nu_1+\nu_2$; then $w(\nu)=2(l-k)$, and $(A_{-2})^k(A_2)^l\subset N^\nu$, as $A_{-2}A_{-2}A_2\subset N_2$. If
$k\le l$, then $c\in(A_2)^{l-k}$ by \cref{pf:vwin-he:fg}, and $b_0^kc=f^kg^kc\in(A_{-2})^k(A_2)^l$; if $k>l$, then
$c\in(A_{-2})^{k-l}$, and $b_0^lc\in(A_{-2})^k(A_2)^l$.

(ii) Let $K_\mu$ be the set of $c\in A$ such that $b_0^nct^\mu$ is an element of $J'$ for some $n$; it is an additive
subgroup of $A$, as $J'$ is an ideal and $b_0\in R_0$. Each of the generators $F_i^2$, $G_i$, $E_i^2$, $E'_i$, $Z_i$ of
$I_2$ is an element $y\in A_{w(\nu)}$ with $yt^\nu\in J'$ for some $\nu\in\{0,e_0,e_1,e_2\}$: for $F_i^2$ ($i<r$) and
$G_i$ take $\nu=e_0$, for $E_i^2$ and $E'_i$ take $\nu=e_1$, and for $F_r^2$ take $\nu=e_2$ (these are entries of $ts$),
and for $Z_i$ take $\nu=0$. Let $c\in A$ and $\nu'=\mu-\nu$, which is $\ge0$ as $\mu\ge e_0+e_1+e_2$. Since $y$ is
homogeneous, the component of degree $w(\mu)$ of $cy$ is $c_{w(\nu')}y$; by (i), $b_0^nc_{w(\nu')}t^{\nu'}\in R$ for
some $n$, and $b_0^nc_{w(\nu')}y\,t^\mu=(b_0^nc_{w(\nu')}t^{\nu'})(yt^\nu)\in J'$. So $(cy)_{w(\mu)}\in K_\mu$, and as
$K_\mu$ is a group, $z_{w(\mu)}\in K_\mu$ for every $z\in I_2$.

(iii) Let $\varphi:N^\mu\to(R/J')_\mu$ map $h$ to the class of $ht^\mu$, and $\psi:N^\mu\to(A/I_2)_{w(\mu)}$ map $h$
to its class; $\varphi$ is surjective. If $\psi(h)=0$, then $h=h_{w(\mu)}\in I_2$, so $b_0^n\varphi(h)=0$ for some $n$
by (ii), and $\sigma(b_0^n)\ne0$. Now \cref{pf:vwin-he:loclen}\ref{pf:vwin-he:loclen-surj} gives the assertion.
\end{proof}

\begin{proof}[Proof of \cref{pf:vwin-he:ann}]
By \cref{pf:vwin-he:stdR} and the remark following it, $(B,R)$ is a standard setting, and $R$ is Noetherian, being
generated over the Noetherian ring $B$ by the finitely generated modules $R_{e_i}$. The quotient $M=R/J'$ is a graded
module, and its pieces $M_d=(R/J')_d$ are finitely generated $B$-modules by \ref{pf:vwin-he:E-gr}. We use three
results of the formal development:
\begin{enumerate}[label=(\roman*)]
\item\label{pf:vwin-he:ann-kill} if $M$ is a graded module over a standard setting $(B,R)$ and
$\mathfrak a,\mathfrak b\subset B$ are ideals with $(\mathfrak aR\cdot\vwinheIrr)^n\subset\vwinheAnn(M)+\mathfrak bR$,
then $\mathfrak a^nM_d\subset\mathfrak bM_d$ for large $d$ (\lref{UpperTheory.V.Alg.M.kill_holds});
\item\label{pf:vwin-he:ann-zi} every prime ideal $\mathfrak p\subset R$ with $\vwinheIrr\not\subset\mathfrak p$ which
contains $\vwinhesecs(ts)$ and the $Z_i$ contains the image of $B\cap I$ (\lref{UpperTheory.V.Alg.M.zi_holds});
\item\label{pf:vwin-he:ann-zii} under the hypotheses of \cref{pf:vwin-he:ann}, every prime ideal $\mathfrak p$ as in
\ref{pf:vwin-he:ann-zi} with $I_DR\subset\mathfrak p$ contains $\mathfrak m_BR$
(\lref{UpperTheory.V.Alg.M.zii_holds}).
\end{enumerate}
Let $\mathfrak a,\mathfrak b\subset B$ be ideals such that every prime ideal $\mathfrak p\subset R$ with
$\vwinheIrr\not\subset\mathfrak p$ and $J'+\mathfrak bR\subset\mathfrak p$ contains $\mathfrak aR$. Then every prime
ideal containing $J'+\mathfrak bR$ contains $\mathfrak aR\cdot\vwinheIrr$; so $\mathfrak aR\cdot\vwinheIrr$ lies in the
radical of $J'+\mathfrak bR$, and as $R$ is Noetherian, $(\mathfrak aR\cdot\vwinheIrr)^n\subset J'+\mathfrak bR$ for
some $n$. Since $J'\subset\vwinheAnn(M)$, \ref{pf:vwin-he:ann-kill} gives $\mathfrak a^nM_d\subset\mathfrak bM_d$ for
large $d$. By \ref{pf:vwin-he:ann-zi} and \ref{pf:vwin-he:ann-zii} this applies to $(\mathfrak a,\mathfrak b)=(B\cap
I,0)$ and to $(\mathfrak a,\mathfrak b)=(\mathfrak m_B,I_D)$: there are $n$ and $n'$ with $(B\cap I)^nM_d=0$ and
$\mathfrak m_B^{n'}M_d\subset I_DM_d$ for large $d$. Fix such a $d$.

Let $L=\vwinheAnn_B(M_d)$, so that $(B\cap I)^n\subset L$. The ring $B/(B\cap I)$ is the image of $B$ in $A/I$, a
finitely generated $\OO$-module by (M2), and so is its quotient $B/(L+B\cap I)$. The ideal
$\mathfrak n=(L+B\cap I)/L$ of the Noetherian ring $B/L$ satisfies $\mathfrak n^n=0$, and each
$\mathfrak n^i/\mathfrak n^{i+1}$ is a finitely generated module over $B/(L+B\cap I)$, hence over $\OO$; so $B/L$ is a
finitely generated $\OO$-module.

Let $\sigma:B\to K$ be an $\OO$-algebra homomorphism. If $\sigma$ vanishes on $I_D$, then
\cref{pf:vwin-he:loclen}\ref{pf:vwin-he:loclen-zero} with $(\mathfrak a,\mathfrak b)=(\mathfrak m_B,I_D)$ and $b=\varpi$
(which lies in $\mathfrak m_B$, as $\OO\to B$ is local, and satisfies $\sigma(\varpi)\ne0$) gives
$\vwinheloclen_\sigma(M_d)=0$. If $\sigma(b)\ne0$ for some $b\in B\cap I$, the same lemma with
$(\mathfrak a,\mathfrak b)=(B\cap I,0)$ gives $\vwinheloclen_\sigma(M_d)=0$.
\end{proof}

\subsubsection{Counting}\label{sec:pf:vwin-he:pf-count}

\begin{proof}[Proof of \cref{pf:vwin-he:num}]
No variable has a weight outside $\{-2,-1,1,2\}$, as $n_{r,s,r',s',0}$ vanishes there. For $i\in\tau$ let
$w(i)=(\max(-\mathrm{wt}(i),0),\max(\mathrm{wt}(i),0))$, which is $(1,0)$, $(2,0)$, $(0,1)$ or $(0,2)$ according as
$\mathrm{wt}(i)$ is $-1$, $-2$, $1$ or $2$; thus $(|\nu|_-,|\nu|_+)=\sum_i\nu_iw(i)$. For $U\subset\tau$ and $p\in\Z^2$
let $N_U(p)$ be the (finite) number of $\nu\in\N^\tau$ supported on $U$ with $\sum_i\nu_iw(i)=p$; thus
$\mathrm{cnt}(a,b)=N_\tau(a,b)$. For $v\in\Z^2$ let $S_v$ be the shift operator $(S_vh)(p)=h(p-v)$ on functions
$h:\Z^2\to\Z$; these operators commute, and $S_{v+v'}=S_vS_{v'}$. Write $X=S_{(1,0)}$ and $Y=S_{(0,1)}$, and let
$\lambda:\Z^3\to\Z^2$, $\lambda(d)=(2d_0+4d_2,2d_1+2d_2)$. As $\lambda$ is additive with $\lambda(e_0)=(2,0)$,
$\lambda(e_1)=(0,2)$ and $\lambda(e_2)=(4,2)$, induction on the length of $L$ gives
\begin{equation}\label{pf:vwin-he:eq-num1}
 \Delta_L(N_\tau\circ\lambda)(d)=(\mathcal DN_\tau)(\lambda(d)),\qquad
 \mathcal D=(1-X^2)^{r-1+s}(1-Y^2)^{r'+s'}(1-X^4Y^2).
\end{equation}
If $i\notin U$, splitting according to whether $\nu_i=0$ gives $N_{U\cup\{i\}}=N_U+S_{w(i)}N_{U\cup\{i\}}$, that is,
$(1-S_{w(i)})N_{U\cup\{i\}}=N_U$; hence $\prod_{i\in U}(1-S_{w(i)})\,N_\tau=N_{\tau\setminus U}$ for every
$U\subset\tau$. Choose variables $i_0$, $i_1$, $i_2$ of weights $-1$, $1$, $2$; they exist as $r\ge1$. Let
$C=(1-X)^{r-1}(1-X^2)^s(1-Y)^{r'}(1-Y^2)^{s'}$. Counting the variables of each weight,
\[
 \prod_{i\ne i_0,i_1}(1-S_{w(i)})=C(1-Y^2),\qquad\prod_{i\ne i_1,i_2}(1-S_{w(i)})=C(1-X),
\]
so that $C(1-Y^2)N_\tau=N_{\{i_0,i_1\}}$ and $C(1-X)N_\tau=N_{\{i_1,i_2\}}$. Using $1-X^2=(1+X)(1-X)$,
$1-Y^2=(1+Y)(1-Y)$ and $1-X^4Y^2=(1-Y^2)+Y^2(1+X^2)(1+X)(1-X)$, we obtain
\[
 \mathcal DN_\tau=(1+X)^{r-1}(1+Y)^{r'}\bigl(N_{\{i_0,i_1\}}+Y^2(1+X^2)(1+X)N_{\{i_1,i_2\}}\bigr).
\]
Now $N_{\{i_0,i_1\}}(a,b)=1$ for $a,b\ge0$, the only monomial being $X_{i_0}^aX_{i_1}^b$, and
$N_{\{i_1,i_2\}}(a,b)=0$ for $a\ge1$, as these two variables have negative weight $0$. A shift of a function which is
eventually equal to a constant $\kappa$ is eventually equal to $\kappa$; so the bracket is eventually equal to $1$, and
$\mathcal DN_\tau$ is eventually equal to $2^{r-1+r'}$. Since $\lambda(d)$ is large when $d$ is large,
\eqref{pf:vwin-he:eq-num1} proves the claim.
\end{proof}

\begin{proof}[Proof of \cref{pf:vwin-he:la}]
For every valid list $ts=((\delta_1,s_1),\dots,(\delta_n,s_n))$ we have, for large $d$,
\begin{equation}\label{pf:vwin-he:eq-small}
 \chi(ts,M/\beta M)(d)=\sum_{S\subset\{1,\dots,n\}}(-1)^{|S|}h\Bigl(d-\sum_{i\in S}\delta_i\Bigr)=\Delta_{\deg ts}h(d),
\end{equation}
by \ref{pf:vwin-he:E-small} applied to the graded module $M/\beta M$ and by expanding the iterated difference. A list
admitted for an ideal remains admitted when it is permuted, or when an entry $s$ is replaced by a divisor of $s$, as
$\vwinhesecs$ does not change, resp.\ grows.

\emph{Step 1: for $y\in R_\delta$, a valid list $ts$ and $c\in\N$ such that $((\delta,\beta^cy))\cdot ts$ is admitted
for $\vwinheAnn(M)$, we have for large $d$}
\begin{equation}\label{pf:vwin-he:eq-peel}
 \chi\bigl(((\delta,\beta^cy))\cdot ts,M\bigr)(d)=c\,\Delta_{\deg ts}h(d)+\chi\bigl(((\delta,y))\cdot ts,M\bigr)(d).
\end{equation}
For $c=0$ there is nothing to prove. For $c\ge1$ write $\beta^cy=\beta\cdot\beta^{c-1}y$; the list
$((0,\beta),(\delta,\beta^{c-1}y))\cdot ts$ is valid, so by \ref{pf:vwin-he:E-prod}, \ref{pf:vwin-he:E-nzd} and
\eqref{pf:vwin-he:eq-small}, for large $d$,
\[
 \chi\bigl(((\delta,\beta^cy))\cdot ts,M\bigr)(d)=\chi(ts,M/\beta M)(d)+\chi\bigl(((\delta,\beta^{c-1}y))\cdot ts,M\bigr)(d)
\]
with $\chi(ts,M/\beta M)(d)=\Delta_{\deg ts}h(d)$; as $((\delta,\beta^{c-1}y))\cdot ts$ is again admitted, induction on
$c$ gives \eqref{pf:vwin-he:eq-peel}.

\emph{Step 2: induction on $n$.} For $n=0$ the left side of the assertion vanishes, and $\chi(us,M)(d)\ge0$ for large
$d$ by \ref{pf:vwin-he:E-pos}. Let $n\ge1$ and $L=((\delta,\beta^{c_i}y_i))_{2\le i\le n}$. By
\eqref{pf:vwin-he:eq-peel} for $y_1$, $c_1$ and the list $L\cdot us$, whose list of degrees is
$((\delta)^{n-1},\deg us)$,
\[
 \chi(ts'',M)(d)=c_1\,\Delta_{((\delta)^{n-1},\deg us)}h(d)+\chi\bigl(((\delta,y_1))\cdot L\cdot us,M\bigr)(d)
\]
for large $d$. The list $((\delta,y_1))\cdot L\cdot us$ is admitted for $\vwinheAnn(M)$, so by \ref{pf:vwin-he:E-perm}
it may be replaced by its permutation $L\cdot((\delta,y_1))\cdot us$, which is admitted as well. The induction
hypothesis, applied to $y_2,\dots,y_n$, $c_2,\dots,c_n$ and the valid list $((\delta,y_1))\cdot us$, shows that
$\chi(L\cdot((\delta,y_1))\cdot us,M)(d)$ is at least $(\sum_{i\ge2}c_i)\,\Delta_{((\delta)^{n-1},\deg us)}h(d)$ for
large $d$ (for $n=1$ it says that this Euler characteristic is non-negative). Adding the two contributions gives the
assertion.
\end{proof}

\subsection{Descent along cyclotomic layers}\label{sec:pf:vwin-he:descent}

In this subsection $\ell$ is a prime, $q=\ell$ if $\ell$ is odd and $q=4$ if $\ell=2$, $L_0$ is a number field
containing a primitive $q$-th root of unity, $\Omega$ is an algebraic closure of $L_0$, $G=\Gal(\Omega/L_0)$ and
$\vwinhecyc:G\to\Zl^\times$ is the cyclotomic character. Let $c_0$ be the largest integer $c$ such that $L_0$ contains a
primitive $\ell^c$-th root of unity, and for $n\ge0$ let $G(n)=\{g\in G:\vwinhecyc(g)\equiv1\pmod{\ell^n}\}$. A
\emph{layer of degree $\ell^m$} is an extension $L_1/L_0$ of degree $\ell^m$ generated by the $\ell^k$-th roots of
unity for some $k$. For a number field $F$ and $e\in\Z$, let $\mathrm{shaLen}(F,\ell,e)$ be the $\Zl$-length of the
group $\Sha_D(e)$ of \cref{def:window:sha}, for any algebraic closure of $F$ and any $\ell$-datum $D$; it depends only
on $(F,\ell,e)$ (\lref{UpperTheory.V.LDataCardLeFinrank.shaLen_eq_shaInv}: the length of $\Sha_D(e)$ equals that of
the group of classes of $\Gal(\overline F/F)$ which vanish on the inertia groups of the places not above $\ell$ and on
the decomposition groups of the places above $\ell$, for a fixed algebraic closure $\overline F$).

\begin{definition}[The Iwasawa module and its dual]\label{pf:vwin-he:dinf}
Let $U\subset G$ be a closed normal subgroup, $U_\infty=U\cap\ker\vwinhecyc$, and $N\ge1$. Let $\vwinheD_N(U)$ be the
$\Zl$-module of locally constant homomorphisms $\varphi:U_\infty\to\vwinheQZ$ which vanish on $U_\infty\cap I_A$ for
every valuation ring $A$ of $\Omega$, and on $U_\infty\cap D_A$ for every valuation ring $A$ whose maximal ideal
contains $N$. The group $G$ acts on $\vwinheD_N(U)$ by $(g\varphi)(u)=\varphi(g^{-1}ug)$, and for $e\in\Z$ we write
$g\ast_e\varphi=\vwinhecyc(g)^{-e}\,g\varphi$. Let $\vwinheXX_N(U)=\Hom_{\Zl}(\vwinheD_N(U),\vwinheQZ)$, with the
action $(g\ast_ex)(\varphi)=\vwinhecyc(g)^e\,x(g^{-1}\varphi)$; thus $(g\ast_ex)(\varphi)=x(g^{-1}\ast_e\varphi)$. For a
subgroup $V\subset G$ let $\vwinheD_N(U)^{V,e}$ be the submodule of elements fixed by $V$ for the action $\ast_e$. For
an endomorphism $T$ of a $\Zl$-module $M$ let $\vwinhecoker(T)=\len_{\Zl}(M/TM)\in\N\cup\{\infty\}$.
\formallabel{UpperTheory.V.Kum.Dinf}
\formallabel{UpperTheory.V.Tower.xTw}
\formallabel{UpperTheory.V.Tower.xTw_apply_apply}
\end{definition}

Thus $\vwinheD_\ell(G)$ is the group of characters of order a power of $\ell$ of the Galois group of the maximal
abelian extension of $L_0(\mu_{\ell^\infty})$ which is unramified at all places and in which the places above $\ell$
split completely, i.e.\ the Pontryagin dual of the module $X'_\infty$ of \cref{lit:window:fw}.

\begin{theorem}[Descent formula]\label{pf:vwin-he:step0}
Let $\gamma\in G$ be such that $\vwinhecyc(G)$ is contained in the closure of the subgroup generated by
$\vwinhecyc(\gamma)$. Let $e\ne0$, $m\ge0$, and let $L_1/L_0$ be a layer of degree $\ell^m$. Then
\[
 \mathrm{shaLen}(L_1,\ell,e)=\vwinhecoker\bigl(t^{\ell^m}-1\bigr),
\]
where $t$ is the action of $\gamma$ on $\vwinheXX_\ell(G)$ by $\ast_e$.
\formallabel{UpperTheory.V.Tower.step0}
\end{theorem}

This is the identity behind the tower bounds of \cref{lem:window:tower}: the Tate--Shafarevich group of a twist
over the layer of degree $\ell^m$ is measured by the coinvariants of $\gamma^{\ell^m}$ on the twisted module
$\vwinheXX_\ell(G)$. The proof occupies the rest of this subsection.

\subsubsection{Layers}

\begin{lemma}\label{pf:vwin-he:layers}
\begin{enumerate}[label=(\alph*)]
\item\label{pf:vwin-he:layers-a} $L_0$ is totally imaginary, $c_0\ge1$ ($c_0\ge2$ if $\ell=2$), and
$\vwinhecyc(G)=1+\ell^{c_0}\Zl$. For $c\ge0$, $L_0$ contains a primitive $\ell^c$-th root of unity if and only if
$c\le c_0$, if and only if $G(c)=G$.
\item\label{pf:vwin-he:layers-b} If $L_1/L_0$ is a layer of degree $\ell^m$ and $L_1$ is embedded in $\Omega$ over
$L_0$, then $\Gal(\Omega/L_1)=G(c_0+m)$, and $[G:G(c_0+m)]=\ell^m$.
\item\label{pf:vwin-he:layers-c} If $\vwinhecyc(G)$ is contained in the closure of the group generated by
$\vwinhecyc(\gamma)$, then $\vwinhecyc(\gamma)\in(1+\ell^{c_0}\Zl)\smallsetminus(1+\ell^{c_0+1}\Zl)$, $\gamma\in G(1)$, and
for every $m\ge0$, $\gamma^{\ell^m}\in G(c_0+m)$ and $\vwinhecyc(G(c_0+m))$ is contained in the closure of the group
generated by $\vwinhecyc(\gamma^{\ell^m})$.
\end{enumerate}
\formallabel{UpperTheory.V.Tower.muQ}
\formallabel{UpperTheory.V.Tower.cycRange}
\formallabel{UpperTheory.V.Tower.layerUniq}
\formallabel{UpperTheory.V.Tower.layerGen}
\end{lemma}

\begin{proof}
A primitive $q$-th root of unity is not real, so $L_0$ is totally imaginary, and $c_0\ge1$, resp.\ $c_0\ge2$. The
group $\vwinhecyc(G)\cong\Gal(L_0(\mu_{\ell^\infty})/L_0)$ is an open, hence infinite, closed subgroup of
$1+q\Zl$, and $L_0$ contains the $\ell^c$-th roots of unity if and only if $\vwinhecyc(G)\subset1+\ell^c\Zl$. For
$c\ge1$ ($c\ge2$ if $\ell=2$) the $\ell$-adic logarithm is an isomorphism $1+\ell^c\Zl\cong\ell^c\Zl$ of topological
groups \cite[Ch.~II, \S5]{Neukirch1999}, so the closed subgroups of $1+q\Zl$ are $1$ and the groups $1+\ell^c\Zl$;
this gives $\vwinhecyc(G)=1+\ell^{c_0}\Zl$ and \ref{pf:vwin-he:layers-a}. If $L_1=L_0(\mu_{\ell^k})$, then
$\Gal(\Omega/L_1)=G(k)$ and $[G:G(k)]=[1+\ell^{c_0}\Zl:1+\ell^{\max(k,c_0)}\Zl]=\ell^{\max(k,c_0)-c_0}$; as this is
$\ell^m$, $\max(k,c_0)=c_0+m$ and $G(k)=G(c_0+m)$, which is \ref{pf:vwin-he:layers-b}. Under the logarithm, the
topological generators of $1+\ell^{c_0}\Zl\cong\ell^{c_0}\Zl$ are the elements of valuation exactly $c_0$, and
raising to the power $\ell^m$ multiplies by $\ell^m$; this gives \ref{pf:vwin-he:layers-c}.
\end{proof}

\subsubsection{Duality}

For a $\Zl$-module $M$ put $M^\vee=\Hom_{\Zl}(M,\vwinheQZ)$, and for a homomorphism $T:M\to M'$ let
$T^\vee:M'^\vee\to M^\vee$, $x\mapsto x\circ T$.

\begin{lemma}\label{pf:vwin-he:dual}
\begin{enumerate}[label=(\alph*)]
\item\label{pf:vwin-he:dual-a} $\vwinheQZ$ is an injective $\Zl$-module, and for every $0\ne m\in M$ there is
$x\in M^\vee$ with $x(m)\ne0$.
\item\label{pf:vwin-he:dual-b} For $T:M\to M'$, an element $x\in M^\vee$ lies in the image of $T^\vee$ if and only if
$x$ vanishes on $\ker T$; and for every submodule $P\subset M$ the restriction $M^\vee\to P^\vee$ is surjective.
\item\label{pf:vwin-he:dual-c} $\len_{\Zl}M^\vee=\len_{\Zl}M$ in $\N\cup\{\infty\}$.
\item\label{pf:vwin-he:dual-d} For an endomorphism $g$ of $M$, $\vwinhecoker(g^\vee)=\len_{\Zl}\ker g$.
\end{enumerate}
\formallabel{UpperTheory.V.PropA.qzInj}
\formallabel{UpperTheory.V.Tower.qzSep}
\formallabel{UpperTheory.V.Tower.dualKer}
\formallabel{UpperTheory.V.Tower.lenDual}
\formallabel{UpperTheory.V.Tower.cokerDual}
\end{lemma}

\begin{proof}
\ref{pf:vwin-he:dual-a} $\vwinheQZ$ is divisible, hence injective over the principal ideal domain $\Zl$. The cyclic
module $\Zl m\cong\Zl/\mathfrak a$ ($\mathfrak a\ne\Zl$) embeds into $\vwinheQZ$, and the embedding extends to $M$ by
injectivity.
\ref{pf:vwin-he:dual-b} If $x\in M^\vee$ vanishes on $\ker T$, it induces a homomorphism $T(M)\to\vwinheQZ$, which
extends to $y\in M'^\vee$ by \ref{pf:vwin-he:dual-a}; then $x=y\circ T$. The converse is clear, and the surjectivity
of the restriction is injectivity of $\vwinheQZ$.
\ref{pf:vwin-he:dual-c} By \ref{pf:vwin-he:dual-a} and \ref{pf:vwin-he:dual-b}, $M\mapsto M^\vee$ is exact. If $M$ has
finite length, induction on the length reduces the assertion to $M$ simple, i.e.\ $M\cong\F_\ell$, and
$\F_\ell^\vee\cong\F_\ell$. If $M$ has infinite length, it has chains $0=M_0\subsetneq M_1\subsetneq\dots\subsetneq
M_n\subset M$ of every length $n$. The submodules $K_i=\{x\in M^\vee:x(M_i)=0\}$ form a decreasing chain, and
$K_{i-1}\ne K_i$: a non-zero homomorphism $M_i/M_{i-1}\to\vwinheQZ$ (which exists by \ref{pf:vwin-he:dual-a})
extends to an element of $K_{i-1}\smallsetminus K_i$. Hence $M^\vee$ has infinite length.
\ref{pf:vwin-he:dual-d} By \ref{pf:vwin-he:dual-b}, restriction to $\ker g$ induces an isomorphism
$M^\vee/g^\vee(M^\vee)\cong(\ker g)^\vee$; conclude by \ref{pf:vwin-he:dual-c}.
\end{proof}

\subsubsection{The module $\vwinheD$}

\begin{lemma}\label{pf:vwin-he:dinf-basic}
Let $U\subset G$ be a closed normal subgroup and $N\ge1$.
\begin{enumerate}[label=(\alph*)]
\item\label{pf:vwin-he:dinf-tors} Every $\varphi\in\vwinheD_N(U)$ is killed by a power of $\ell$, and $U_\infty$ acts
trivially on $\vwinheD_N(U)$ and on $\vwinheXX_N(U)$, for every twisted action $\ast_e$.
\item\label{pf:vwin-he:dinf-nil} For every $\varphi\in\vwinheD_N(U)$ there is $n$ such that $g\varphi=\varphi$ for all
$g\in U\cap G(n)$.
\item\label{pf:vwin-he:dinf-gen} Let $V\subset U$ be a subgroup and $\gamma\in V$ such that $\vwinhecyc(V)$ is contained
in the closure of the group generated by $\vwinhecyc(\gamma)$. Then for every $e\in\Z$,
$\vwinheD_N(U)^{V,e}$ is the kernel of $\gamma\ast_e-1$ and also the kernel of $\gamma^{-1}\ast_e-1$.
\end{enumerate}
\formallabel{UpperTheory.V.Tower.dinfTors}
\formallabel{UpperTheory.V.Kum.dinfNil}
\formallabel{UpperTheory.V.Tower.shaDGen}
\end{lemma}

\begin{proof}
\ref{pf:vwin-he:dinf-tors} $U_\infty$ is closed in the compact group $G$, so a locally constant $\varphi$ has finite
image, which consists of torsion elements of $\vwinheQZ$. For $u,w\in U_\infty$,
$\varphi(u^{-1}wu)=-\varphi(u)+\varphi(w)+\varphi(u)=\varphi(w)$, and $\vwinhecyc(u)=1$; so $U_\infty$ acts
trivially on $\vwinheD_N(U)$, and then on its dual.

\ref{pf:vwin-he:dinf-nil} The set $\{u\in U_\infty:\varphi(u)=0\}$ is open in $U_\infty$, so there is an open normal
subgroup $W$ of the profinite group $G$ with $\varphi(W\cap U_\infty)=0$. For $w\in W$ and $u\in U_\infty$ we have
$w^{-1}uwu^{-1}\in W\cap U_\infty$, hence $\varphi(w^{-1}uw)=\varphi(u)$: the elements of $W$ fix $\varphi$, and so do
those of $U_\infty$ by \ref{pf:vwin-he:dinf-tors}. Hence every element of the open set $S=WU_\infty$ fixes
$\varphi$. The set $U\smallsetminus S$ is compact, and $\vwinhecyc(U\smallsetminus S)$ is a compact subset of
$\Zl^\times$ not containing $1$, since $U\cap\ker\vwinhecyc=U_\infty\subset S$. So it does not meet $1+\ell^n\Zl$
for large $n$, i.e.\ $U\cap G(n)\subset S$.

\ref{pf:vwin-he:dinf-gen} Invariants under $V$ are invariant under $\gamma$ and $\gamma^{-1}$. Conversely, let
$\gamma\ast_e\varphi=\varphi$ (equivalently $\gamma^{-1}\ast_e\varphi=\varphi$), and $v\in V$. By
\ref{pf:vwin-he:dinf-tors} and \ref{pf:vwin-he:dinf-nil} choose $r$ with $\ell^r\varphi=0$ and $n\ge r$ such that
$U\cap G(n)$ fixes $\varphi$; then every $w\in U\cap G(n)$ satisfies $w\ast_e\varphi=\vwinhecyc(w)^{-e}\varphi=\varphi$,
since $\vwinhecyc(w)^{-e}\equiv1\pmod{\ell^r}$. Since $\vwinhecyc(v)$ lies in the closure of the powers of
$\vwinhecyc(\gamma)$, there is $j\in\Z$ with $\vwinhecyc(\gamma^{-j}v)\in1+\ell^n\Zl$, so
$\gamma^{-j}v\in U\cap G(n)$ and $v\ast_e\varphi=\gamma^j\ast_e(\gamma^{-j}v)\ast_e\varphi=\varphi$.
\end{proof}

\subsubsection{Restriction to the cyclotomic kernel}

Let $F$ be a number field with algebraic closure $\Omega$, $G_F=\Gal(\Omega/F)$, $\vwinhecyc$ its cyclotomic
character and $U_\infty(F)=\ker\vwinhecyc$. We say that $F$ is \emph{procyclic at $\ell$} if $\ell$ is odd or
$\vwinhecyc(G_F)\subset1+4\Z_2$; by \cref{pf:vwin-he:layers}\ref{pf:vwin-he:layers-a} this holds if $F$ contains a
primitive $q$-th root of unity. Let $D$ be an $\ell$-datum for $F$ (\cref{def:window:sha}), with set $S=S_\ell$ of
places above $\ell$ and group $G_{F,S}$, and let $e\in\Z$. For a continuous cocycle $z$ of $G_{F,S}$ with values in
$\vwinheQZ(-e)$ (on which $G_{F,S}$ acts through $\vwinhecyc^{-e}$), the function $u\mapsto z(\bar u)$ on
$U_\infty(F)$, $\bar u$ the image of $u$ in $G_{F,S}$, only depends on the class of $z$, since $U_\infty(F)$ acts
trivially; this defines the \emph{restriction} $\mathrm{res}:H^1(G_{F,S},\vwinheQZ(-e))\to\mathrm{Map}(U_\infty(F),\vwinheQZ)$.
With the notation of \cref{pf:vwin-he:dinf} for $F$ in place of $L_0$, let $\mathcal E_e(F)$ be the set of
$\varphi\in\vwinheD_\ell(G_F)$ such that $\varphi(gug^{-1})=\vwinhecyc(g)^{-e}\varphi(u)$ for all $g\in G_F$ and
$u\in U_\infty(F)$.

\begin{lemma}[Generator criterion]\label{pf:vwin-he:cob}
Let $\Gamma$ be a compact group, $\omega:\Gamma\to\Zl^\times$ a continuous homomorphism, $M$ a discrete abelian group
with a continuous action of $\Gamma$ on which $\ker\omega$ acts trivially, and $\gamma\in\Gamma$ such that
$\omega(\Gamma)$ is contained in the closure of the group generated by $\omega(\gamma)$ and $m\mapsto\gamma m-m$ is
surjective on $M$. Then every continuous cocycle $z:\Gamma\to M$ vanishing on $\ker\omega$ is a coboundary.
\formallabel{UpperTheory.V.Kum.cobOfGen}
\end{lemma}

\begin{proof}
Choose $m_0$ with $\gamma m_0-m_0=z(\gamma)$ and put $z'(g)=z(g)-(gm_0-m_0)$. Then $z'$ is a continuous cocycle
vanishing at $\gamma$ and on $\ker\omega$ (on which coboundaries vanish). For $n\in\ker\omega$,
$z'(gn)=z'(g)+gz'(n)=z'(g)$, so $z'=\bar z\circ\omega$ for a function $\bar z$ on $\omega(\Gamma)$, which is
continuous because $\omega:\Gamma\to\omega(\Gamma)$ is a continuous surjection from a compact space, hence a quotient
map. By the cocycle relation $z'(\gamma^k)=0$ for all $k\in\Z$. Since $\omega(\Gamma)$ is compact, it is the closure of
the group generated by $\omega(\gamma)$, and $\bar z$, being locally constant and zero on a dense subset, vanishes.
Hence $z$ is the coboundary of $m_0$.
\end{proof}

\begin{lemma}[Injectivity]\label{pf:vwin-he:resinj}
If $e\ne0$ and $F$ is procyclic at $\ell$, then $\mathrm{res}$ is injective.
\formallabel{UpperTheory.V.Kum.resInfInj}
\end{lemma}

\begin{proof}
First, $\vwinhecyc^e$ is not trivial on $G_F$. Otherwise let $n=|e|$, $k=n[F:\Q]+1$ and $\zeta\in\Omega$ a primitive
$\ell^k$-th root of unity; for $g\in G_F$, $g\zeta=\zeta^{\vwinhecyc(g)}$ with $\vwinhecyc(g)$ an $n$-th root of unity
in $\Zl$, of which there are at most $n$. So $[F(\zeta):F]\le n$ and $[\Q(\zeta):\Q]\le n[F:\Q]=k-1$, whereas
$[\Q(\zeta):\Q]=\ell^{k-1}(\ell-1)\ge k$.

The group $\vwinhecyc(G_F)$ is a closed subgroup of $\Zl^\times$, contained in $1+4\Z_2$ if $\ell=2$; so it contains
an element $\vwinhecyc(\gamma)$ generating a dense subgroup (\lref{UpperTheory.V.Kum.existsGen}: such closed
subgroups of $\Zl^\times$ are topologically generated by one element). Then $\vwinhecyc(\gamma)^{-e}\ne1$, for
otherwise $\vwinhecyc^e$ would be trivial on $G_F$; hence $\gamma$ acts on $\vwinheQZ(-e)$ as multiplication by
$\vwinhecyc(\gamma)^{-e}$, and $m\mapsto\gamma m-m$ is multiplication by a non-zero element of $\Zl$, which is
surjective on the divisible group $\vwinheQZ$. If $\mathrm{res}(z)=0$, the cocycle $z$ vanishes on the image of
$U_\infty(F)$, which is the kernel of $\vwinhecyc$ on $G_{F,S}$. By \cref{pf:vwin-he:cob} for $\Gamma=G_{F,S}$, $z$ is a
coboundary.
\end{proof}

\begin{lemma}\label{pf:vwin-he:vanish}
Let $M$ be an abelian group in which every element is killed by a power of $\ell$, $\varphi:U_\infty(F)\to M$ a
locally constant homomorphism, and $\gamma\in G_F$ with $\vwinhecyc(\gamma)$ of infinite order. Then $\varphi$ vanishes
on $U_\infty(F)\cap C$, where $C$ is the closure of the group generated by $\gamma$.
\formallabel{UpperTheory.V.Kum.vanishCyc}
\end{lemma}

\begin{proof}
Let $x\in U_\infty(F)\cap C$. There is an open normal subgroup $W\subset G_F$ with $\varphi(W\cap U_\infty(F))=0$.
Choose $d\ge1$ with $\gamma^d\in W$ and write $d=\ell^ad'$ with $\ell\nmid d'$. As $\vwinhecyc(\gamma)$ has infinite
order, there is $j$ such that $\vwinhecyc(\gamma)^k\equiv1\pmod{\ell^j}$ implies $\ell^a\mid k$ (use the logarithm
on $1+\ell^2\Zl$ as in the proof of \cref{pf:vwin-he:layers}). Since $x\in C$, there is $k\in\Z$ with
$x^{-1}\gamma^k\in W\cap G_F(j)$, where $G_F(j)=\{g:\vwinhecyc(g)\equiv1\pmod{\ell^j}\}$; as $\vwinhecyc(x)=1$,
$\vwinhecyc(\gamma)^k\equiv1\pmod{\ell^j}$ and $\ell^a\mid k$. The elements $x$ and $\gamma$ commute, so
$x^{d'}=(\gamma^k)^{d'}\bigl((x^{-1}\gamma^k)^{d'}\bigr)^{-1}$, where $(\gamma^k)^{d'}=(\gamma^d)^{k/\ell^a}\in W$ and
$x^{-1}\gamma^k\in W$. Hence $x^{d'}\in W\cap U_\infty(F)$ and $d'\varphi(x)=\varphi(x^{d'})=0$; since $\varphi(x)$ has
$\ell$-power order and $\ell\nmid d'$, $\varphi(x)=0$.
\end{proof}

\begin{lemma}[Extension of equivariant homomorphisms]\label{pf:vwin-he:ext}
Assume that $F$ is procyclic at $\ell$. Let $M$ be a discrete abelian group in which every element is killed by a
power of $\ell$, with a continuous action of $G_F$ which is trivial on $U_\infty(F)$, and let $\varphi:U_\infty(F)\to M$
be a locally constant homomorphism with $\varphi(gug^{-1})=g\varphi(u)$ for $g\in G_F$, $u\in U_\infty(F)$. Then there is
a continuous cocycle $z:G_F\to M$ with $z|_{U_\infty(F)}=\varphi$.
\formallabel{UpperTheory.V.Kum.extCocycle}
\end{lemma}

\begin{proof}
As in the proof of \cref{pf:vwin-he:resinj} choose $\gamma\in G_F$ such that $\vwinhecyc(\gamma)$ generates a dense
subgroup of $\vwinhecyc(G_F)$; since $\vwinhecyc(G_F)\cong\Gal(F(\mu_{\ell^\infty})/F)$ is infinite,
$\vwinhecyc(\gamma)$ has infinite order. Let $C$ be the closure of the group generated by $\gamma$. Since $\vwinhecyc(C)$
is compact and contains the powers of $\vwinhecyc(\gamma)$, $\vwinhecyc(C)=\vwinhecyc(G_F)$, so $G_F=C\cdot U_\infty(F)$.
Put $z(cu)=c\varphi(u)$ for $c\in C$, $u\in U_\infty(F)$. This is well defined: if $cu=c'u'$, then
$c'^{-1}c=u'u^{-1}\in C\cap U_\infty(F)$, on which $\varphi$ vanishes by \cref{pf:vwin-he:vanish}; hence
$\varphi(u')=\varphi(u)$, and $c\varphi(u)=c'\varphi(u)$ since $c'^{-1}c$ acts trivially. The multiplication map
$C\times U_\infty(F)\to G_F$ is a continuous surjection of compact spaces, hence a quotient map, and
$(c,u)\mapsto c\varphi(u)$ is continuous; so $z$ is continuous. For $c,c'\in C$ and $u,u'\in U_\infty(F)$,
$cu\cdot c'u'=cc'\cdot(c'^{-1}uc')u'$, and using $\varphi(u)=c'\varphi(c'^{-1}uc')$ and the triviality of the action of
$u$,
\[
 z(cuc'u')=cc'\varphi(c'^{-1}uc')+cc'\varphi(u')=c\varphi(u)+cu\cdot c'\varphi(u')=z(cu)+cu\,z(c'u').
\qedhere
\]
\end{proof}

\begin{lemma}[Local triviality above $\ell$]\label{pf:vwin-he:deccob}
Assume that $e\ne0$ and that $F$ is procyclic at $\ell$. Let $A$ be a valuation ring of $\Omega$ above a place $v$ of $F$
above $\ell$, and $z:G_F\to\vwinheQZ(-e)$ a continuous cocycle vanishing on $D_A\cap U_\infty(F)$. Then the restriction
of $z$ to $D_A$ is a coboundary.
\formallabel{UpperTheory.V.Kum.decCob}
\end{lemma}

\begin{proof}
$D_A$ is closed in $G_F$, hence compact, and $\vwinhecyc(D_A)$ is open in $\Zl^\times$
(\lref{UpperTheory.V.LDataCardLeFinrank.isOpen_decImage}; indeed $D_A$ is the Galois group of the completion
$F_v$, and $[F_v(\mu_{\ell^n}):F_v]\ge\ell^{n-1}(\ell-1)/[F_v:\Ql]$, so $\vwinhecyc(D_A)$ is a closed subgroup of finite
index). As in the proof of \cref{pf:vwin-he:resinj}, it has a topological generator $\vwinhecyc(g_0)$, $g_0\in D_A$,
which has infinite order since $\vwinhecyc(D_A)$ is infinite; so $\vwinhecyc(g_0)^{-e}\ne1$ and $m\mapsto g_0m-m$ is
surjective. Apply \cref{pf:vwin-he:cob} to $\Gamma=D_A$ and $\omega=\vwinhecyc|_{D_A}$.
\end{proof}

\begin{lemma}[Image of the restriction]\label{pf:vwin-he:resimage}
If $e\ne0$ and $F$ is procyclic at $\ell$, then $\mathrm{res}(\Sha_D(e))=\mathcal E_e(F)$.
\formallabel{UpperTheory.V.Kum.resInfMem}
\formallabel{UpperTheory.V.Kum.resInfOnto}
\end{lemma}

\begin{proof}
Let $z$ be a cocycle of $G_{F,S}$ representing a class of $\Sha_D(e)$, viewed as a cocycle of $G_F$, and
$\varphi=\mathrm{res}(z)$. Since $U_\infty(F)$ acts trivially, $\varphi$ is a homomorphism, and it is locally constant.
It vanishes on $I_A$ for $A$ not above $\ell$, since $z$ is inflated from $G_{F,S}$. For $A$ above $\ell$, the
restriction of $z$ to $D_A$ is a coboundary $g\mapsto gm-m$ (for the places $A_v$ of the datum by definition, and for
the other places above $\ell$ because they are conjugate to these), which vanishes on $D_A\cap U_\infty(F)$; this
includes $I_A\subset D_A$. Finally, for $g\in G_F$ and $u\in U_\infty(F)$ the cocycle relation and the triviality of
the action of $u$ give $z(gug^{-1})=z(g)+g\,z(u)-gug^{-1}z(g)=g\,z(u)=\vwinhecyc(g)^{-e}z(u)$. So
$\varphi\in\mathcal E_e(F)$.

Conversely let $\varphi\in\mathcal E_e(F)$. By \cref{pf:vwin-he:ext} (with $M=\vwinheQZ(-e)$) there is a continuous
cocycle $z$ of $G_F$ extending $\varphi$. For $A$ not above $\ell$, the inertia group $I_A$ is contained in
$U_\infty(F)$, as $\vwinhecyc$ is unramified outside $\ell$, and $\varphi$ vanishes on it; so $z$ vanishes on the closed
normal subgroup generated by these inertia groups, which acts trivially on $\vwinheQZ(-e)$, and $z$ factors through a
cocycle of $G_{F,S}$. For $A$ above $\ell$, $z$ vanishes on $D_A\cap U_\infty(F)$, so its restriction to $D_A$ is a
coboundary by \cref{pf:vwin-he:deccob}. Hence the class of $z$ lies in $\Sha_D(e)$, and $\mathrm{res}(z)=\varphi$.
\end{proof}

\begin{lemma}[Descent]\label{pf:vwin-he:desc}
Let $L_1\subset\Omega$ be a finite extension of $L_0$, procyclic at $\ell$, with Galois group
$G_{L_1}=\Gal(\Omega/L_1)\subset G$, let $U\subset G$ be a closed normal subgroup with
$U\cap\ker\vwinhecyc=G_{L_1}\cap\ker\vwinhecyc$, let $e\ne0$, and let $D$ be an $\ell$-datum for $L_1$. Then restriction
to $U_\infty$ induces an isomorphism of $\Zl$-modules $\Sha_D(e)\cong\vwinheD_\ell(U)^{G_{L_1},e}$.
\formallabel{UpperTheory.V.Tower.desc}
\formallabel{UpperTheory.V.Tower.dinfTransport}
\end{lemma}

\begin{proof}
The groups $U_\infty=U\cap\ker\vwinhecyc$ and $U_\infty(L_1)=G_{L_1}\cap\ker\vwinhecyc$ coincide, and for a valuation ring
$A$ of $\Omega$ the inertia and decomposition groups of $A$ in $G_{L_1}$ are those in $G$ intersected with $G_{L_1}$;
since $U_\infty\subset G_{L_1}$, the conditions defining $\vwinheD_\ell(U)$ (for $L_0$) and $\vwinheD_\ell(G_{L_1})$
(for $L_1$) agree, and the two modules are the same set of functions on $U_\infty$. For $\varphi$ in this module,
$g\ast_e\varphi=\varphi$ means $\vwinhecyc(g)^{-e}\varphi(g^{-1}ug)=\varphi(u)$ for all $u$; for all $g\in G_{L_1}$ this
is the condition $\varphi(gug^{-1})=\vwinhecyc(g)^{-e}\varphi(u)$ defining $\mathcal E_e(L_1)$. By
\cref{pf:vwin-he:resinj,pf:vwin-he:resimage} for $F=L_1$, restriction is an isomorphism
$\Sha_D(e)\to\mathcal E_e(L_1)=\vwinheD_\ell(U)^{G_{L_1},e}$.
\end{proof}

\begin{proof}[Proof of \cref{pf:vwin-he:step0}]
Embed $L_1$ into $\Omega$ over $L_0$ and choose an $\ell$-datum $D$ for $L_1$. As $L_1\supset L_0$ contains a primitive
$q$-th root of unity, $L_1$ is procyclic at $\ell$. By \cref{pf:vwin-he:layers}\ref{pf:vwin-he:layers-b},
$V:=G_{L_1}=G(c_0+m)$, and $V\cap\ker\vwinhecyc=G\cap\ker\vwinhecyc$. \Cref{pf:vwin-he:desc} with $U=G$ gives
$\Sha_D(e)\cong\vwinheD^{V,e}$, where $\vwinheD=\vwinheD_\ell(G)$. By
\cref{pf:vwin-he:layers}\ref{pf:vwin-he:layers-c}, $\gamma'=\gamma^{\ell^m}$ lies in $V$ and $\vwinhecyc(V)$ is contained in the
closure of the group generated by $\vwinhecyc(\gamma')$; by
\cref{pf:vwin-he:dinf-basic}\ref{pf:vwin-he:dinf-gen}, $\vwinheD^{V,e}=\ker g$ for the endomorphism
$g=\gamma'^{-1}\ast_e-1$ of $\vwinheD$. Hence $\mathrm{shaLen}(L_1,\ell,e)=\len_{\Zl}\Sha_D(e)=\len_{\Zl}\ker g$. On
$\vwinheXX_\ell(G)=\vwinheD^\vee$ we have $t^{\ell^m}=\gamma'\ast_e$, and for $x\in\vwinheD^\vee$ and
$\varphi\in\vwinheD$,
\[
 \bigl((t^{\ell^m}-1)x\bigr)(\varphi)=x(\gamma'^{-1}\ast_e\varphi)-x(\varphi)=x(g\varphi),
\]
i.e.\ $t^{\ell^m}-1=g^\vee$. By \cref{pf:vwin-he:dual}\ref{pf:vwin-he:dual-d},
$\vwinhecoker(t^{\ell^m}-1)=\len_{\Zl}\ker g$.
\end{proof}

\begin{remark}\label{pf:vwin-he:rem-descent}
(1) The formula holds for all $e\ne0$, of either sign, and for $\ell=2$ (under the assumption $\mu_4\subset L_0$, which
is what makes $\Gal(L_0(\mu_{2^\infty})/L_0)$ procyclic). No finiteness is assumed: both sides may be infinite.
(2) The module $\vwinheXX_\ell(G)$ is the full $\Zl$-dual of $\vwinheD_\ell(G)$, i.e.\ the double dual of $X'_\infty$;
no topology on $\vwinheXX_\ell(G)$ is needed, since only lengths of cokernels enter.
(3) An element $\gamma$ as in the theorem exists (\lref{UpperTheory.V.Tower.genGal}: the absolute Galois group of
a number field which is procyclic at $\ell$ contains an element whose cyclotomic image topologically generates the
image of the cyclotomic character).
\end{remark}

